% ============================================================
% NUMERICAL ANALYSIS — Lecture Notes (L2/L3)
% English version — Compilation: xelatex -shell-escape
% ============================================================
\documentclass[12pt,a4paper,twoside]{book}

% --- Encoding and language (XeLaTeX) ---
\usepackage{fontspec}
\usepackage{polyglossia}
\setmainlanguage{english}

% --- Page layout ---
\usepackage[margin=2.5cm]{geometry}
\usepackage{fancyhdr}
\pagestyle{fancy}
\fancyhead[LE]{\leftmark}
\fancyhead[RO]{\rightmark}
\fancyfoot[C]{\thepage}

% --- Mathematics ---
\usepackage{amsmath,amssymb,amsthm,mathtools}
\usepackage{bm}

% --- Graphics ---
\usepackage{tikz,pgfplots}
\pgfplotsset{compat=1.18}
\usetikzlibrary{patterns,arrows.meta,calc,positioning,decorations.pathreplacing,shapes}

% --- Code ---
\usepackage{minted}
\setminted{fontsize=\small,frame=lines,framesep=2mm,breaklines=true}
\setminted[python]{python3=true}

% --- Algorithms ---
\usepackage[ruled,vlined]{algorithm2e}
\SetKwInput{KwInput}{Input}
\SetKwInput{KwOutput}{Output}

% --- Tables ---
\usepackage{booktabs,array,multirow,longtable}

% --- Colors and boxes ---
\usepackage[dvipsnames]{xcolor}
\usepackage{tcolorbox}
\tcbuselibrary{theorems,breakable,skins}

% --- Links ---
\usepackage{hyperref}
\hypersetup{colorlinks=true,linkcolor=NavyBlue,citecolor=ForestGreen,urlcolor=Maroon}

% --- Bibliography and index ---
\usepackage[numbers]{natbib}
\usepackage{makeidx}
\makeindex

% ============================================================
% Theorem environments
% ============================================================
\theoremstyle{definition}
\newtheorem{definition}{Definition}[chapter]
\newtheorem{example}[definition]{Example}
\newtheorem{exercise}{Exercise}[chapter]

\theoremstyle{plain}
\newtheorem{theorem}[definition]{Theorem}
\newtheorem{proposition}[definition]{Proposition}
\newtheorem{lemma}[definition]{Lemma}
\newtheorem{corollary}[definition]{Corollary}

\theoremstyle{remark}
\newtheorem{remark}[definition]{Remark}

% ============================================================
% Boxes
% ============================================================
\newtcolorbox{keyformulas}[1][]{colback=blue!5,colframe=NavyBlue,
  title=\textbf{Key Formulas},fonttitle=\bfseries,breakable,#1}
\newtcolorbox{warning}[1][]{colback=red!5,colframe=BrickRed,
  title=\textbf{Warning},fonttitle=\bfseries,breakable,#1}
\newtcolorbox{method}[1][]{colback=green!5,colframe=ForestGreen,
  title=\textbf{Method},fonttitle=\bfseries,breakable,#1}
\newtcolorbox{codeblock}[1][]{colback=gray!5,colframe=gray!80,
  title=\textbf{Implementation},fonttitle=\bfseries,breakable,#1}

% ============================================================
% Commands
% ============================================================
\newcommand{\R}{\mathbb{R}}
\newcommand{\N}{\mathbb{N}}
\newcommand{\C}{\mathbb{C}}
\newcommand{\eps}{\varepsilon}
\newcommand{\fl}{\mathrm{fl}}
\newcommand{\cond}{\mathrm{cond}}
\newcommand{\BigO}{\mathcal{O}}
\DeclareMathOperator{\diag}{diag}
\DeclareMathOperator{\spn}{span}
\DeclareMathOperator{\rang}{rank}

% ============================================================
\begin{document}
\frontmatter
\begin{titlepage}
\begin{tikzpicture}[remember picture, overlay]
  \fill[blue!10!white] (current page.south west) rectangle (current page.north east);
  \fill[blue!80!black, rotate around={-15:(current page.center)}]
    ([xshift=-4cm,yshift=-1cm]current page.center) rectangle ([xshift=4cm,yshift=1.5cm]current page.center);
  \node[anchor=center, text=white, font=\Huge\bfseries, rotate=-15]
    at ([yshift=0.3cm]current page.center) {Numerical Analysis};
  \node[anchor=center, text=orange!80!yellow, font=\Large, rotate=-15]
    at ([yshift=-1cm]current page.center) {Lecture Notes};
  \node[anchor=north, text=blue!70!black, font=\Large\itshape]
    at ([yshift=-4cm]current page.center) {Ya\'e Ulrich Gaba};
  \node[anchor=north, text=blue!50!black, font=\large]
    at ([yshift=-5.5cm]current page.center) {Licence L3 --- 2025--2026};
  \fill[orange!70!yellow] ([yshift=3cm]current page.south west) rectangle ([yshift=3.3cm]current page.south east);
  \node[anchor=south, text=gray, font=\small] at ([yshift=1cm]current page.south) {\today};
\end{tikzpicture}
\end{titlepage}
\tableofcontents
\listofalgorithms
\addcontentsline{toc}{chapter}{List of Algorithms}

% --- Preface ---
\chapter*{Preface}
\addcontentsline{toc}{chapter}{Preface}

Numerical analysis is the art of solving mathematical problems \emph{approximately but with controlled error}. Unlike theoretical analysis, which establishes existence and uniqueness of solutions, numerical analysis is concerned with \emph{actually computing} them and \emph{quantifying the error incurred}.

\medskip
This course covers the fundamentals: floating-point arithmetic, solving linear systems, interpolation, numerical integration, ordinary differential equations, and eigenvalue problems. Each method is accompanied by its error analysis, pseudocode, and implementations in Python and Julia.

\medskip
\textbf{Prerequisites.} Real Analysis I \& II, linear algebra, basic programming.

\medskip
\textbf{References.}
\begin{itemize}
\item \textsc{Quarteroni}, Sacco, Saleri --- \emph{Numerical Mathematics}, Springer.
\item \textsc{S\"uli}, Mayers --- \emph{An Introduction to Numerical Analysis}, Cambridge.
\item \textsc{Trefethen}, Bau --- \emph{Numerical Linear Algebra}, SIAM.
\item \textsc{Higham} --- \emph{Accuracy and Stability of Numerical Algorithms}, SIAM.
\item \textsc{Demailly} --- \emph{Analyse Num\'erique et \'Equations Diff\'erentielles}, EDP Sciences.
\end{itemize}

\mainmatter

% ============================================================
\chapter{Floating-Point Arithmetic, Errors and Stability}
\label{ch:floating}
% ============================================================

\begin{center}
\textit{``Numerical computation is the art of giving the right answer to a slightly different problem.''}
\end{center}

\section{Motivating example}

Let us compute $f(x)=\frac{1-\cos x}{x^2}$ for $x=10^{-8}$ in double precision. The theoretical result is $\approx 0.5$, but Python returns $0.0$! The \textbf{catastrophic cancellation} between $1$ and $\cos x$ eliminates all significant digits. Understanding floating-point arithmetic is essential.

\section{Representation of numbers in a computer}

\begin{definition}[Floating-point number]\index{floating-point number}
In base $\beta$ with $t$ significant digits, a floating-point number is written:
\[
x = \pm m \times \beta^e, \quad m = d_0.d_1d_2\cdots d_{t-1}, \quad d_0\neq 0,
\]
where $e_{\min}\leq e\leq e_{\max}$ and $0\leq d_i<\beta$.
\end{definition}

\begin{definition}[IEEE 754 Standard]\index{IEEE 754}
\begin{center}
\begin{tabular}{lccc}
\toprule
\textbf{Format} & \textbf{Bits} & \textbf{Mantissa} & \textbf{$\eps_{\mathrm{mach}}$} \\
\midrule
Single precision & 32 & 23+1 bits & $\approx 1.19\times 10^{-7}$ \\
Double precision & 64 & 52+1 bits & $\approx 2.22\times 10^{-16}$ \\
\bottomrule
\end{tabular}
\end{center}
\end{definition}

\begin{definition}[Machine epsilon]\index{machine epsilon}
The \textbf{machine epsilon} $\eps_{\mathrm{mach}}$ is the smallest $\eps>0$ such that $\fl(1+\eps)>1$:
\[
\eps_{\mathrm{mach}} = \frac{1}{2}\beta^{1-t}.
\]
For every representable $x\in\R$, $\fl(x)=x(1+\delta)$ with $|\delta|\leq\eps_{\mathrm{mach}}$.
\end{definition}

\section{Rounding errors and propagation}

\begin{definition}[Absolute and relative error]\index{absolute error}\index{relative error}
\begin{align}
\text{Absolute error} &: E_a = |x-\tilde{x}|, \\
\text{Relative error} &: E_r = \frac{|x-\tilde{x}|}{|x|} \quad (x\neq 0).
\end{align}
\end{definition}

\begin{theorem}[Error propagation]
If $\tilde{x}=x(1+\eps_x)$ and $\tilde{y}=y(1+\eps_y)$:
\begin{itemize}
\item $\tilde{x}\cdot\tilde{y}\approx xy(1+\eps_x+\eps_y)$: relative errors add up.
\item $\tilde{x}-\tilde{y}\approx(x-y)+x\eps_x-y\eps_y$: if $x\approx y$, the relative error blows up.
\end{itemize}
\end{theorem}

\begin{warning}
\textbf{Catastrophic cancellation} occurs when subtracting nearly equal numbers. Example: $1-\cos x$ for small $x$. Remedy: use $2\sin^2(x/2)$.
\end{warning}

\section{Conditioning of a problem}

\begin{definition}[Condition number]\index{condition number}
The \textbf{condition number} of a problem $f$ at the point $x$ is:
\[
\kappa(f,x) = \left|\frac{xf'(x)}{f(x)}\right|.
\]
A problem is \textbf{well-conditioned} if $\kappa\sim 1$, \textbf{ill-conditioned} if $\kappa\gg 1$.
\end{definition}

\begin{definition}[Matrix condition number]\index{matrix condition number}
For $Ax=b$:
\[
\cond(A) = \|A\|\cdot\|A^{-1}\|.
\]
The relative error on $x$ is amplified by $\cond(A)$:
\[
\frac{\|\delta x\|}{\|x\|} \leq \cond(A)\frac{\|\delta b\|}{\|b\|}.
\]
\end{definition}

\section{Numerical stability}

\begin{definition}[Stable algorithm]\index{stability}
An algorithm is \textbf{stable} if it gives the exact solution of a \textbf{slightly perturbed} problem. It is \textbf{backward stable} if the perturbation is of order $\eps_{\mathrm{mach}}$.
\end{definition}

\begin{method}
The accuracy of the result depends on two factors:
\begin{enumerate}
\item The \textbf{conditioning} of the problem (inherent, cannot be changed).
\item The \textbf{stability} of the algorithm (the programmer's choice).
\end{enumerate}
Rule: error $\lesssim \cond \times \eps_{\mathrm{mach}}$.
\end{method}

\begin{center}
\begin{tikzpicture}[
  box/.style={rectangle, draw, rounded corners, minimum width=2.5cm, minimum height=1.2cm, text centered, font=\small},
  arr/.style={-{Stealth}, thick}
]
\node[box,fill=blue!10,align=center] (prob) at (0,0) {Problem\\$f(x)$};
\node[box,fill=green!10,align=center] (algo) at (5,0) {Algorithm\\$\tilde{f}(\tilde{x})$};
\node[box,fill=orange!10,align=center] (res) at (10,0) {Result\\$\tilde{y}$};
\draw[arr] (prob) -- node[above]{\small Conditioning $\kappa$} (algo);
\draw[arr] (algo) -- node[above]{\small Stability} (res);
\draw[arr,dashed,red] (prob) to[bend left=30] node[above]{\small Total error $\leq\kappa\cdot\eps$} (res);
\end{tikzpicture}
\end{center}

\section{Implementation}

\begin{codeblock}[title=\textbf{Python}]
\begin{minted}{python}
import numpy as np

# Epsilon machine
eps = np.finfo(float).eps
print(f"eps_mach = {eps}")  # 2.220446049250313e-16

# Cancellation catastrophique
x = 1e-8
f_bad = (1 - np.cos(x)) / x**2
f_good = 2 * np.sin(x/2)**2 / x**2
print(f"Mauvais : {f_bad}")   # 0.0
print(f"Bon     : {f_good}")  # 0.49999999...

# Conditionnement d'une matrice
A = np.array([[1, 1], [1, 1.0001]])
print(f"cond(A) = {np.linalg.cond(A):.1f}")  # ~40000
\end{minted}
\end{codeblock}

\begin{codeblock}[title=\textbf{Julia}]
\begin{minted}{julia}
# Epsilon machine
println("eps_mach = ", eps(Float64))

# Cancellation
x = 1e-8
f_bad = (1 - cos(x)) / x^2
f_good = 2sin(x/2)^2 / x^2
println("Mauvais : $f_bad")
println("Bon     : $f_good")

# Conditionnement
using LinearAlgebra
A = [1.0 1.0; 1.0 1.0001]
println("cond(A) = ", cond(A))
\end{minted}
\end{codeblock}

\section{Exercises}

\begin{exercise}[$\star$ -- Machine epsilon]
Write a program that computes $\eps_{\mathrm{mach}}$ by bisection. Compare with \texttt{np.finfo(float).eps}.
\end{exercise}

\begin{exercise}[$\star$ -- Cancellation]
Propose a stable computation of $\sqrt{x+1}-\sqrt{x}$ for large $x$. Verify numerically.
\end{exercise}

\begin{exercise}[$\star\star$ -- Condition number]
Compute the condition number of the problem $f(x)=\ln x$ at the point $x=1$. Interpret.
\end{exercise}

\begin{exercise}[$\star\star$ -- Hilbert matrix]
Let $H_n$ be the Hilbert matrix ($h_{ij}=1/(i+j-1)$). Compute $\cond(H_n)$ for $n=2,\ldots,15$. What do you observe?
\end{exercise}

\begin{exercise}[$\star\star\star$ -- Project]
Study the numerical stability of evaluating the polynomial $p(x)=\sum a_ix^i$ using Horner's algorithm vs.\ naive evaluation. Compare relative errors for high-degree polynomials.
\end{exercise}

\begin{keyformulas}
\begin{align*}
\fl(x) &= x(1+\delta),\;|\delta|\leq\eps_{\mathrm{mach}} &
\eps_{\mathrm{mach}} &= \tfrac{1}{2}\beta^{1-t} \\
\kappa(f,x) &= |xf'(x)/f(x)| &
\cond(A) &= \|A\|\cdot\|A^{-1}\|
\end{align*}
\end{keyformulas}

% ============================================================
\chapter{Root Finding}
\label{ch:roots}
% ============================================================

\section{Motivating example}
Finding $x$ such that $e^x = 3x$ amounts to finding a root of $f(x)=e^x-3x$. No closed-form solution exists: a numerical method is needed.

\section{Bisection method}

\begin{theorem}[Intermediate Value Theorem]\index{bisection}
If $f\in C([a,b])$ and $f(a)f(b)<0$, then $\exists\,c\in(a,b)$ such that $f(c)=0$.
\end{theorem}

\begin{algorithm}[H]
\caption{Bisection}\label{alg:bisection}
\KwInput{$f$, $a$, $b$, tolerance $\eps$}
\KwOutput{Approximation of the root $c$}
\While{$b-a > \eps$}{
  $c \gets (a+b)/2$\;
  \eIf{$f(a)\cdot f(c) < 0$}{$b \gets c$}{$a \gets c$}
}
\Return{$c$}
\end{algorithm}

\begin{theorem}[Convergence of bisection]
After $n$ iterations: $|c_n-x^*|\leq\frac{b-a}{2^{n+1}}$. The convergence is \textbf{linear} with rate $1/2$.
\end{theorem}

\begin{proof}
At each iteration, the interval is halved. After $n$ iterations, $|I_n|=(b-a)/2^n$ and $|c_n-x^*|\leq|I_n|/2$.
\end{proof}

\section{Newton's method}\index{Newton}

\begin{definition}
Newton's iteration for $f(x)=0$ is:
\[
x_{n+1} = x_n - \frac{f(x_n)}{f'(x_n)}.
\]
\end{definition}

\begin{center}
\begin{tikzpicture}
\begin{axis}[
  width=10cm, height=6cm, axis lines=middle,
  xlabel=$x$, ylabel=$y$, ymin=-2, ymax=5,
  domain=0.5:3, samples=80, title={Newton's method}
]
\addplot[thick,blue] {exp(x) - 3*x};
\addplot[dashed,red] coordinates {(2,{exp(2)-6}) (2-((exp(2)-6)/(exp(2)-3)),0)};
\addplot[only marks,mark=*,red] coordinates {(2,{exp(2)-6})};
\node[red,below] at (axis cs:1.37,0) {\small $x_1$};
\node[red,above] at (axis cs:2,1.4) {\small $(x_0,f(x_0))$};
\end{axis}
\end{tikzpicture}
\end{center}

\begin{algorithm}[H]
\caption{Newton's method}\label{alg:newton}
\KwInput{$f$, $f'$, $x_0$, tolerance $\eps$, $N_{\max}$}
\KwOutput{Approximation $x_n$}
\For{$n=0,1,\ldots,N_{\max}$}{
  $x_{n+1} \gets x_n - f(x_n)/f'(x_n)$\;
  \If{$|x_{n+1}-x_n|<\eps$}{\Return{$x_{n+1}$}}
}
\end{algorithm}

\begin{theorem}[Quadratic convergence of Newton's method]\label{thm:newton-conv}
If $f\in C^2$, $f(x^*)=0$, $f'(x^*)\neq 0$ and $x_0$ is sufficiently close to $x^*$, then:
\[
|e_{n+1}| \leq \frac{M}{2m}|e_n|^2,
\]
where $e_n=x_n-x^*$, $M=\sup|f''|$, $m=\inf|f'|$ in a neighbourhood of $x^*$.
\end{theorem}

\begin{proof}
By Taylor's theorem: $0=f(x^*)=f(x_n)+f'(x_n)(x^*-x_n)+\frac{f''(\xi_n)}{2}(x^*-x_n)^2$. Dividing by $f'(x_n)$:
\[
0 = \frac{f(x_n)}{f'(x_n)}+(x^*-x_n)+\frac{f''(\xi_n)}{2f'(x_n)}e_n^2.
\]
Since $x_{n+1}=x_n-f(x_n)/f'(x_n)$, we get $e_{n+1}=x_{n+1}-x^*=-\frac{f''(\xi_n)}{2f'(x_n)}e_n^2$.
\end{proof}

\section{Secant method}\index{secant method}

\begin{definition}
We replace $f'(x_n)$ by a finite difference:
\[
x_{n+1} = x_n - f(x_n)\frac{x_n-x_{n-1}}{f(x_n)-f(x_{n-1})}.
\]
\end{definition}

\begin{theorem}
The order of convergence of the secant method is the \textbf{golden ratio} $\varphi=(1+\sqrt{5})/2\approx 1.618$.
\end{theorem}

\section{Fixed-point method}\index{fixed point}

\begin{definition}
Find $x=g(x)$. Iteration: $x_{n+1}=g(x_n)$.
\end{definition}

\begin{theorem}[Banach Fixed-Point Theorem]\label{thm:banach}
If $g:[a,b]\to[a,b]$ is a contraction ($|g'(x)|\leq L<1$ on $[a,b]$), then $g$ has a unique fixed point $x^*$ and $x_n\to x^*$ with:
\[
|x_n-x^*|\leq \frac{L^n}{1-L}|x_1-x_0|.
\]
\end{theorem}

\section{Comparison of methods}

\begin{center}
\begin{tabular}{lccc}
\toprule
\textbf{Method} & \textbf{Order} & \textbf{Evaluations of $f$} & \textbf{Robustness} \\
\midrule
Bisection & 1 (linear) & 1/iter. & Very robust \\
Newton & 2 (quadratic) & $f + f'$/iter. & Sensitive to $x_0$ \\
Secant & $\varphi\approx 1.618$ & 1/iter. & Moderate \\
\bottomrule
\end{tabular}
\end{center}

\section{Numerical example}

Root of $f(x)=e^x-3x$ with $x_0=2$ (Newton):
\begin{center}
\begin{tabular}{cccc}
\toprule
$n$ & $x_n$ & $f(x_n)$ & $|x_n-x_{n-1}|$ \\
\midrule
0 & 2.000000 & 1.389056 & --- \\
1 & 1.373418 & 0.225698 & 0.627 \\
2 & 1.524247 & $-0.046$ & 0.151 \\
3 & 1.512127 & $-0.000374$ & 0.012 \\
4 & 1.512135 & $<10^{-9}$ & $8\times10^{-6}$ \\
\bottomrule
\end{tabular}
\end{center}

\section{Implementation}

\begin{codeblock}[title=\textbf{Python}]
\begin{minted}{python}
import numpy as np

def bisection(f, a, b, tol=1e-12):
    while b - a > tol:
        c = (a + b) / 2
        if f(a) * f(c) < 0:
            b = c
        else:
            a = c
    return (a + b) / 2

def newton(f, df, x0, tol=1e-12, maxiter=100):
    x = x0
    for i in range(maxiter):
        dx = f(x) / df(x)
        x -= dx
        if abs(dx) < tol:
            return x, i+1
    return x, maxiter

def secant(f, x0, x1, tol=1e-12, maxiter=100):
    for i in range(maxiter):
        fx0, fx1 = f(x0), f(x1)
        x2 = x1 - fx1 * (x1 - x0) / (fx1 - fx0)
        if abs(x2 - x1) < tol:
            return x2, i+1
        x0, x1 = x1, x2
    return x1, maxiter

# Test
f = lambda x: np.exp(x) - 3*x
df = lambda x: np.exp(x) - 3

print(f"Bissection : {bisection(f, 0, 1):.12f}")
x_newton, n = newton(f, df, 2.0)
print(f"Newton     : {x_newton:.12f} ({n} iterations)")
x_sec, n = secant(f, 0.0, 1.0)
print(f"Secante    : {x_sec:.12f} ({n} iterations)")
\end{minted}
\end{codeblock}

\begin{codeblock}[title=\textbf{Julia}]
\begin{minted}{julia}
function newton(f, df, x0; tol=1e-12, maxiter=100)
    x = x0
    for i in 1:maxiter
        dx = f(x) / df(x)
        x -= dx
        abs(dx) < tol && return x, i
    end
    return x, maxiter
end

f(x) = exp(x) - 3x
df(x) = exp(x) - 3
x, n = newton(f, df, 2.0)
println("Newton: x = $x ($n iterations)")
\end{minted}
\end{codeblock}

\section{Exercises}

\begin{exercise}[$\star$]
Find $\sqrt{2}$ by applying Newton's method to $f(x)=x^2-2$. How many iterations are needed for 15 decimal digits?
\end{exercise}

\begin{exercise}[$\star\star$]
Show that Newton's method applied to $f(x)=x^2$ converges linearly (double root). Generalise: what happens when $f'(x^*)=0$?
\end{exercise}

\begin{exercise}[$\star\star$]
Prove that the order of the secant method is $\varphi$ by setting $e_{n+1}\approx C\,e_n^p$ and solving $p^2-p-1=0$.
\end{exercise}

\begin{exercise}[$\star\star\star$ -- Project]
Implement Brent's method (combination of bisection + secant). Compare with the simpler methods on 10 test functions.
\end{exercise}

\begin{keyformulas}
\begin{align*}
\text{Bisection} &: |e_n|\leq\frac{b-a}{2^{n+1}} \\
\text{Newton} &: x_{n+1}=x_n-\frac{f(x_n)}{f'(x_n)},\quad |e_{n+1}|\leq C|e_n|^2 \\
\text{Secant} &: x_{n+1}=x_n-f(x_n)\frac{x_n-x_{n-1}}{f(x_n)-f(x_{n-1})},\quad \text{order }\varphi
\end{align*}
\end{keyformulas}

% ============================================================
\chapter{Direct Methods for Linear Systems}
\label{ch:direct}
% ============================================================

\section{Motivating example}
An electrical circuit with $n$ nodes produces a system $Ax=b$ of $n$ equations. For $n=1000$, a systematic and efficient method is needed.

\section{Gaussian elimination and LU factorisation}

\begin{definition}[LU factorisation]\index{LU factorisation}\index{LU}
Find $L$ lower triangular (with $\ell_{ii}=1$) and $U$ upper triangular such that $A=LU$. Then $Ax=b$ is solved by $Ly=b$ (forward substitution) followed by $Ux=y$ (back substitution).
\end{definition}

\begin{algorithm}[H]
\caption{LU factorisation (without pivoting)}\label{alg:lu}
\KwInput{Matrix $A\in\R^{n\times n}$}
\KwOutput{$L$, $U$ such that $A=LU$}
\For{$k=1,\ldots,n-1$}{
  \For{$i=k+1,\ldots,n$}{
    $\ell_{ik} \gets a_{ik}/a_{kk}$\;
    \For{$j=k+1,\ldots,n$}{
      $a_{ij} \gets a_{ij}-\ell_{ik}\cdot a_{kj}$\;
    }
  }
}
$U \gets$ upper triangular part of $A$\;
\end{algorithm}

\begin{theorem}[Existence of LU]
If all leading principal submatrices of $A$ are invertible, then the LU factorisation exists and is unique.
\end{theorem}

\begin{theorem}[Cost]
The LU factorisation costs $\frac{2n^3}{3}+\BigO(n^2)$ floating-point operations.
\end{theorem}

\begin{proof}
Step $k$ performs $(n-k)^2$ multiplications and subtractions. Total: $\sum_{k=1}^{n-1}2(n-k)^2=2\sum_{j=1}^{n-1}j^2=\frac{2n^3}{3}+\BigO(n^2)$.
\end{proof}

\section{Partial pivoting}\index{partial pivoting}

\begin{warning}
Without pivoting, the algorithm can be unstable (division by a small pivot). \textbf{Partial pivoting} permutes rows to maximise $|a_{kk}|$.
\end{warning}

One obtains $PA=LU$ where $P$ is a permutation matrix.

\begin{theorem}[Stability]
With partial pivoting, Gaussian elimination is \textbf{backward stable}:
\[
(A+\Delta A)\hat{x}=b, \quad \|\Delta A\|\leq g(n)\,\eps_{\mathrm{mach}}\,\|A\|,
\]
where $g(n)$ is the growth factor (bounded by $2^{n-1}$ in theory, rarely exceeded in practice).
\end{theorem}

\section{Cholesky factorisation}\index{Cholesky}

\begin{theorem}[Cholesky factorisation]
If $A$ is \textbf{symmetric positive definite} (SPD), then there exists a unique lower triangular matrix $L$ with positive diagonal entries such that $A=LL^\top$.
\end{theorem}

\begin{proof}
By induction on $n$. Write $A=\begin{pmatrix}a_{11}&\mathbf{v}^\top\\\mathbf{v}&A'\end{pmatrix}$. Since $A$ is SPD, $a_{11}>0$. Set $\ell_{11}=\sqrt{a_{11}}$, $\mathbf{l}=\mathbf{v}/\ell_{11}$. Then $A'- \mathbf{l}\mathbf{l}^\top$ is SPD of size $(n-1)\times(n-1)$ and we apply the induction hypothesis.
\end{proof}

\begin{algorithm}[H]
\caption{Cholesky factorisation}\label{alg:cholesky}
\KwInput{$A$ SPD, $n\times n$}
\KwOutput{$L$ such that $A=LL^\top$}
\For{$j=1,\ldots,n$}{
  $\ell_{jj}\gets\sqrt{a_{jj}-\sum_{k=1}^{j-1}\ell_{jk}^2}$\;
  \For{$i=j+1,\ldots,n$}{
    $\ell_{ij}\gets\frac{1}{\ell_{jj}}\left(a_{ij}-\sum_{k=1}^{j-1}\ell_{ik}\ell_{jk}\right)$\;
  }
}
\end{algorithm}

\begin{theorem}[Cost of Cholesky]
$\frac{n^3}{3}+\BigO(n^2)$: half the cost of LU.
\end{theorem}

\section{Numerical example}

$A=\begin{pmatrix}4&2\\2&5\end{pmatrix}$, $b=\begin{pmatrix}8\\9\end{pmatrix}$. Cholesky: $L=\begin{pmatrix}2&0\\1&2\end{pmatrix}$.

Forward substitution $Ly=b$: $y_1=4$, $y_2=(9-1\cdot4)/2=2.5$. Back substitution $L^\top x=y$: $x_2=2.5/2=1.25$, $x_1=(4-1\cdot 1.25)/2=1.375$.

\section{Implementation}

\begin{codeblock}[title=\textbf{Python}]
\begin{minted}{python}
import numpy as np
from scipy.linalg import lu, cholesky, solve_triangular

# LU
A = np.array([[2., 1, 1], [4, 3, 3], [8, 7, 9]])
b = np.array([4., 10, 24])
P, L, U = lu(A)
y = solve_triangular(L, P @ b, lower=True)
x = solve_triangular(U, y)
print(f"LU: x = {x}")

# Cholesky
A_spd = np.array([[4., 2], [2, 5]])
b_spd = np.array([8., 9])
L_chol = cholesky(A_spd, lower=True)
y = solve_triangular(L_chol, b_spd, lower=True)
x = solve_triangular(L_chol.T, y)
print(f"Cholesky: x = {x}")
print(f"Cout LU (n=100): ~{2*100**3//3:.0f} flops")
\end{minted}
\end{codeblock}

\begin{codeblock}[title=\textbf{Julia}]
\begin{minted}{julia}
using LinearAlgebra
A = [2.0 1 1; 4 3 3; 8 7 9]
b = [4.0, 10, 24]
F = lu(A)
x = F \ b
println("LU: x = $x")

A_spd = [4.0 2; 2 5]
b_spd = [8.0, 9]
C = cholesky(A_spd)
x = C \ b_spd
println("Cholesky: x = $x")
\end{minted}
\end{codeblock}

\section{Exercises}

\begin{exercise}[$\star$]
Factorise $A=\begin{pmatrix}1&2&3\\2&8&14\\3&14&34\end{pmatrix}$ into $LU$ by hand. Verify.
\end{exercise}

\begin{exercise}[$\star\star$]
Show that if $A$ is SPD, all pivots in Gaussian elimination are positive (no pivoting needed).
\end{exercise}

\begin{exercise}[$\star\star$]
Compare the computation times of LU vs.\ Cholesky for $n=100,500,1000,2000$ on random SPD matrices. Verify the ratio $\sim 2$.
\end{exercise}

\begin{exercise}[$\star\star\star$ -- Project]
Implement the banded LU factorisation for tridiagonal matrices (Thomas algorithm). Apply to the discretisation of $-u''=f$.
\end{exercise}

\begin{keyformulas}
\begin{align*}
A &= LU \quad (\text{cost } \tfrac{2n^3}{3}) \\
A &= LL^\top \quad (\text{Cholesky, cost } \tfrac{n^3}{3}) \\
PA &= LU \quad (\text{partial pivoting}) \\
\cond(A) &= \|A\|\cdot\|A^{-1}\|, \quad \frac{\|\delta x\|}{\|x\|}\leq\cond(A)\frac{\|\delta b\|}{\|b\|}
\end{align*}
\end{keyformulas}

% ============================================================
\chapter{Iterative Methods for Linear Systems}
\label{ch:iterative}
% ============================================================

\section{Motivating example}
In finite element methods, the matrix $A$ is sparse ($\BigO(n)$ nonzero entries out of $n^2$). LU costs $\BigO(n^3)$, which is unacceptable for $n=10^6$. Iterative methods exploit the sparse structure: each iteration costs $\BigO(n)$.

\section{General principle}

\begin{definition}[Splitting method]
We write $A=M-N$ with $M$ easily invertible. The iteration is:
\[
Mx^{(k+1)} = Nx^{(k)}+b, \quad\text{i.e.}\quad x^{(k+1)}=Gx^{(k)}+c,
\]
where $G=M^{-1}N$ is the \textbf{iteration matrix} and $c=M^{-1}b$.
\end{definition}

\begin{theorem}[Convergence]\label{thm:iter-conv}
The method converges for every $x^{(0)}$ if and only if $\rho(G)<1$ (spectral radius).
\end{theorem}

\begin{proof}
$e^{(k)}=x^{(k)}-x^*=G^k e^{(0)}$. Now $\|G^k\|\to 0 \iff \rho(G)<1$.
\end{proof}

\section{Jacobi method}\index{Jacobi}

\begin{definition}
$A=D-E-F$ (diagonal, lower triangular, upper triangular). $M=D$:
\[
x_i^{(k+1)} = \frac{1}{a_{ii}}\left(b_i-\sum_{j\neq i}a_{ij}x_j^{(k)}\right).
\]
Iteration matrix: $G_J=D^{-1}(E+F)$.
\end{definition}

\section{Gauss--Seidel method}\index{Gauss--Seidel}

\begin{definition}
$M=D-E$:
\[
x_i^{(k+1)} = \frac{1}{a_{ii}}\left(b_i-\sum_{j<i}a_{ij}x_j^{(k+1)}-\sum_{j>i}a_{ij}x_j^{(k)}\right).
\]
Iteration matrix: $G_{GS}=(D-E)^{-1}F$.
\end{definition}

\begin{theorem}[Convergence for SPD matrices]
If $A$ is SPD, Gauss--Seidel converges.
\end{theorem}

\begin{theorem}[Convergence for diagonally dominant matrices]
If $A$ is strictly diagonally dominant ($|a_{ii}|>\sum_{j\neq i}|a_{ij}|$), then both Jacobi and Gauss--Seidel converge.
\end{theorem}

\section{Successive over-relaxation (SOR)}\index{SOR}

\begin{definition}
We introduce a parameter $\omega\in(0,2)$:
\[
x_i^{(k+1)} = (1-\omega)x_i^{(k)} + \frac{\omega}{a_{ii}}\left(b_i-\sum_{j<i}a_{ij}x_j^{(k+1)}-\sum_{j>i}a_{ij}x_j^{(k)}\right).
\]
$\omega=1$: Gauss--Seidel. $\omega>1$: over-relaxation. $\omega<1$: under-relaxation.
\end{definition}

\begin{theorem}
SOR converges if and only if $0<\omega<2$. The optimal $\omega$ depends on $\rho(G_J)$.
\end{theorem}

\section{Conjugate gradient method}\index{conjugate gradient}

\begin{definition}
For $A$ SPD, solving $Ax=b$ is equivalent to minimising $\phi(x)=\frac{1}{2}x^\top Ax-b^\top x$.
\end{definition}

\begin{algorithm}[H]
\caption{Conjugate gradient}\label{alg:cg}
\KwInput{$A$ SPD, $b$, $x_0$, tolerance $\eps$}
\KwOutput{$x\approx A^{-1}b$}
$r_0\gets b-Ax_0$, $p_0\gets r_0$\;
\For{$k=0,1,2,\ldots$}{
  $\alpha_k\gets\frac{r_k^\top r_k}{p_k^\top Ap_k}$\;
  $x_{k+1}\gets x_k+\alpha_k p_k$\;
  $r_{k+1}\gets r_k-\alpha_k Ap_k$\;
  \If{$\|r_{k+1}\|<\eps$}{\Return{$x_{k+1}$}}
  $\beta_k\gets\frac{r_{k+1}^\top r_{k+1}}{r_k^\top r_k}$\;
  $p_{k+1}\gets r_{k+1}+\beta_k p_k$\;
}
\end{algorithm}

\begin{theorem}[Convergence of the conjugate gradient method]
In exact arithmetic, CG converges in at most $n$ iterations. The error satisfies:
\[
\|e_k\|_A \leq 2\left(\frac{\sqrt{\kappa}-1}{\sqrt{\kappa}+1}\right)^k\|e_0\|_A, \quad \kappa=\cond_2(A).
\]
\end{theorem}

\section{Implementation}

\begin{codeblock}[title=\textbf{Python}]
\begin{minted}{python}
import numpy as np

def jacobi(A, b, x0, tol=1e-10, maxiter=1000):
    D = np.diag(np.diag(A))
    R = A - D
    x = x0.copy()
    for k in range(maxiter):
        x_new = np.linalg.solve(D, b - R @ x)
        if np.linalg.norm(x_new - x) < tol:
            return x_new, k+1
        x = x_new
    return x, maxiter

def conjugate_gradient(A, b, x0, tol=1e-10, maxiter=1000):
    x = x0.copy()
    r = b - A @ x
    p = r.copy()
    rs_old = r @ r
    for k in range(maxiter):
        Ap = A @ p
        alpha = rs_old / (p @ Ap)
        x += alpha * p
        r -= alpha * Ap
        rs_new = r @ r
        if np.sqrt(rs_new) < tol:
            return x, k+1
        p = r + (rs_new / rs_old) * p
        rs_old = rs_new
    return x, maxiter

# Test : matrice tridiagonale
n = 100
A = np.diag(2*np.ones(n)) - np.diag(np.ones(n-1), 1) - np.diag(np.ones(n-1), -1)
b = np.ones(n)
x0 = np.zeros(n)

x_j, nj = jacobi(A, b, x0)
x_cg, ncg = conjugate_gradient(A, b, x0)
print(f"Jacobi: {nj} iterations")
print(f"CG:     {ncg} iterations")
\end{minted}
\end{codeblock}

\begin{codeblock}[title=\textbf{Julia}]
\begin{minted}{julia}
function cg(A, b, x0; tol=1e-10, maxiter=1000)
    x = copy(x0)
    r = b - A * x
    p = copy(r)
    rs = dot(r, r)
    for k in 1:maxiter
        Ap = A * p
        alpha = rs / dot(p, Ap)
        x .+= alpha .* p
        r .-= alpha .* Ap
        rs_new = dot(r, r)
        sqrt(rs_new) < tol && return x, k
        p .= r .+ (rs_new / rs) .* p
        rs = rs_new
    end
    return x, maxiter
end
\end{minted}
\end{codeblock}

\section{Exercises}

\begin{exercise}[$\star$]
Apply 3 iterations of Jacobi and Gauss--Seidel to the system $\begin{pmatrix}4&1\\1&3\end{pmatrix}x=\begin{pmatrix}1\\2\end{pmatrix}$, $x^{(0)}=(0,0)^\top$.
\end{exercise}

\begin{exercise}[$\star\star$]
Show that for a tridiagonal matrix $T_n$ of size $n$, $\rho(G_J)=\cos(\pi/(n+1))$. Deduce the number of Jacobi iterations needed for a precision $\eps$.
\end{exercise}

\begin{exercise}[$\star\star\star$ -- Project]
Compare Jacobi, Gauss--Seidel, SOR and CG for the 2D discretisation of the Laplacian $-\Delta u=f$ on $[0,1]^2$. Plot the number of iterations as a function of $n$.
\end{exercise}

\begin{keyformulas}
\begin{align*}
\text{Jacobi} &: x^{(k+1)}=D^{-1}(b-(A-D)x^{(k)}) \\
\text{Gauss--Seidel} &: G=(D-E)^{-1}F \\
\text{CG} &: \|e_k\|_A\leq 2\left(\frac{\sqrt{\kappa}-1}{\sqrt{\kappa}+1}\right)^k\|e_0\|_A
\end{align*}
\end{keyformulas}

% ============================================================
\chapter{Polynomial Interpolation}
\label{ch:interpolation}
% ============================================================

\section{Motivating example}
We have temperature measurements at 5 time instants. We want to estimate the temperature at an intermediate time. Polynomial interpolation constructs a polynomial passing exactly through the measured points.

\section{The interpolation problem}

\begin{theorem}[Existence and uniqueness]
Given $n+1$ distinct points $(x_0,y_0),\ldots,(x_n,y_n)$, there exists a \textbf{unique} polynomial $p_n\in\mathcal{P}_n$ such that $p_n(x_i)=y_i$ for all $i$.
\end{theorem}

\begin{proof}
The Vandermonde matrix $V_{ij}=x_i^j$ is invertible since $\det(V)=\prod_{i>j}(x_i-x_j)\neq 0$.
\end{proof}

\section{Lagrange form}\index{Lagrange}

\begin{definition}
\[
p_n(x) = \sum_{i=0}^n y_i\,L_i(x), \quad L_i(x)=\prod_{\substack{j=0\\j\neq i}}^n\frac{x-x_j}{x_i-x_j}.
\]
The $L_i$ are the \textbf{Lagrange basis polynomials}: $L_i(x_j)=\delta_{ij}$.
\end{definition}

\section{Newton form}\index{Newton!interpolation}

\begin{definition}[Divided differences]
\[
f[x_i]=f(x_i), \quad f[x_i,\ldots,x_{i+k}]=\frac{f[x_{i+1},\ldots,x_{i+k}]-f[x_i,\ldots,x_{i+k-1}]}{x_{i+k}-x_i}.
\]
\end{definition}

\begin{definition}[Newton form]
\[
p_n(x) = f[x_0]+f[x_0,x_1](x-x_0)+\cdots+f[x_0,\ldots,x_n]\prod_{j=0}^{n-1}(x-x_j).
\]
Evaluation by the generalised Horner algorithm: $\BigO(n)$ operations.
\end{definition}

\section{Interpolation error}

\begin{theorem}[Interpolation error]\label{thm:interp-error}
If $f\in C^{n+1}([a,b])$, then for every $x\in[a,b]$:
\[
f(x)-p_n(x) = \frac{f^{(n+1)}(\xi_x)}{(n+1)!}\prod_{i=0}^n(x-x_i),
\]
where $\xi_x\in(a,b)$ depends on $x$.
\end{theorem}

\begin{proof}
Set $w(t)=\prod(t-x_i)$ and $\varphi(t)=f(t)-p_n(t)-\lambda w(t)$ with $\lambda$ chosen so that $\varphi(x)=0$ (for a fixed $x\neq x_i$). Then $\varphi$ vanishes at $n+2$ points. By Rolle's theorem applied $n+1$ times, $\varphi^{(n+1)}$ vanishes at some point $\xi$. Since $\varphi^{(n+1)}=f^{(n+1)}-\lambda(n+1)!$, we get $\lambda=f^{(n+1)}(\xi)/(n+1)!$.
\end{proof}

\begin{warning}
The \textbf{Runge phenomenon}: for $f(x)=1/(1+25x^2)$ on $[-1,1]$ with equidistant nodes, the interpolation error \emph{diverges} as $n\to\infty$ near the endpoints.
\end{warning}

\begin{center}
\begin{tikzpicture}
\begin{axis}[
  width=10cm, height=6cm, domain=-1:1, samples=100,
  xlabel=$x$, ylabel=$y$, ymin=-0.5, ymax=1.5,
  title={Runge phenomenon}
]
\addplot[thick,blue] {1/(1+25*x^2)};
\addlegendentry{$f(x)$}
\end{axis}
\end{tikzpicture}
\end{center}

\section{Chebyshev nodes}\index{Chebyshev!nodes}

\begin{definition}
The Chebyshev nodes on $[-1,1]$:
\[
x_k = \cos\left(\frac{2k+1}{2(n+1)}\pi\right), \quad k=0,\ldots,n.
\]
They minimise $\max_{x\in[-1,1]}|\prod(x-x_i)|$.
\end{definition}

\begin{theorem}
With Chebyshev nodes: $\max|w(x)|\leq 2^{-n}$, which avoids the Runge phenomenon for analytic functions.
\end{theorem}

\section{Cubic splines}\index{spline}

\begin{definition}[Cubic interpolating spline]
$s\in C^2([a,b])$, $s|_{[x_i,x_{i+1}]}\in\mathcal{P}_3$, $s(x_i)=y_i$. Boundary conditions:
\begin{itemize}
\item \textbf{Natural}: $s''(x_0)=s''(x_n)=0$.
\item \textbf{Clamped}: $s'(x_0)=f'(x_0)$, $s'(x_n)=f'(x_n)$.
\end{itemize}
\end{definition}

\begin{theorem}[Cubic spline error]
If $f\in C^4$ and $h=\max(x_{i+1}-x_i)$:
\[
\|f-s\|_\infty \leq \frac{5}{384}h^4\|f^{(4)}\|_\infty.
\]
\end{theorem}

\section{Implementation}

\begin{codeblock}[title=\textbf{Python}]
\begin{minted}{python}
import numpy as np
from scipy.interpolate import lagrange, CubicSpline
import matplotlib.pyplot as plt

# Interpolation de Lagrange
x_nodes = np.array([-1, -0.5, 0, 0.5, 1])
f = lambda x: 1 / (1 + 25*x**2)
y_nodes = f(x_nodes)
p = lagrange(x_nodes, y_nodes)

x_fine = np.linspace(-1, 1, 300)
plt.figure(figsize=(10, 5))
plt.plot(x_fine, f(x_fine), 'b-', label='f(x)', lw=2)
plt.plot(x_fine, p(x_fine), 'r--', label=f'Lagrange deg {len(x_nodes)-1}')
plt.plot(x_nodes, y_nodes, 'ko', markersize=8)
plt.legend(); plt.title("Interpolation de Lagrange")
plt.savefig("ch05_lagrange.pdf"); plt.show()

# Noeuds de Tchebychev
n = 15
cheb = np.cos((2*np.arange(n+1)+1)/(2*(n+1))*np.pi)
y_cheb = f(cheb)
p_cheb = lagrange(cheb, y_cheb)
print(f"Erreur equidist (n=15): {np.max(np.abs(f(x_fine)-lagrange(np.linspace(-1,1,16), f(np.linspace(-1,1,16)))(x_fine))):.4f}")
print(f"Erreur Tchebychev (n=15): {np.max(np.abs(f(x_fine)-p_cheb(x_fine))):.6f}")

# Spline cubique
cs = CubicSpline(x_nodes, y_nodes)
print(f"Spline en 0.25: {cs(0.25):.6f}, exact: {f(0.25):.6f}")
\end{minted}
\end{codeblock}

\begin{codeblock}[title=\textbf{Julia}]
\begin{minted}{julia}
using Interpolations

f(x) = 1 / (1 + 25x^2)
nodes = [-1.0, -0.5, 0, 0.5, 1.0]
vals = f.(nodes)

# Differences divisees (Newton)
function divided_diff(x, y)
    n = length(x)
    F = copy(y)
    for j in 2:n, i in n:-1:j
        F[i] = (F[i] - F[i-1]) / (x[i] - x[i-j+1])
    end
    return F
end
println("Coefs Newton: ", divided_diff(nodes, vals))
\end{minted}
\end{codeblock}

\section{Exercises}

\begin{exercise}[$\star$]
Construct the Lagrange polynomial passing through $(0,1)$, $(1,3)$, $(2,7)$ and verify.
\end{exercise}

\begin{exercise}[$\star\star$]
Prove Theorem~\ref{thm:interp-error} in detail.
\end{exercise}

\begin{exercise}[$\star\star$]
Compare the interpolation error with equidistant nodes vs.\ Chebyshev nodes for $f(x)=1/(1+25x^2)$, $n=5,10,15,20$.
\end{exercise}

\begin{exercise}[$\star\star\star$ -- Project]
Implement Hermite interpolation (data for both $f$ and $f'$). Apply to the construction of B\'ezier curves.
\end{exercise}

\begin{keyformulas}
\begin{align*}
L_i(x) &= \prod_{j\neq i}\frac{x-x_j}{x_i-x_j} \\
f(x)-p_n(x) &= \frac{f^{(n+1)}(\xi)}{(n+1)!}\prod(x-x_i) \\
x_k^{\text{Cheb}} &= \cos\!\left(\frac{2k+1}{2(n+1)}\pi\right) \\
\|f-s\|_\infty &\leq \frac{5}{384}h^4\|f^{(4)}\|_\infty
\end{align*}
\end{keyformulas}

%!TEX root = ../main.tex
% Chapter 6 — Approximation: Least Squares and Chebyshev

\chapter{Approximation --- Least Squares and Chebyshev}
\label{ch:approximation}


\section{Introduction}

In many practical situations, we have a set of data points and wish to find a function that \emph{best represents} them. Unlike interpolation, which requires the approximating function to pass exactly through every data point, approximation seeks the function that is \emph{closest} in some well-defined sense.

\begin{definition}[Approximation problem]
\label{def:approximation}
\index{approximation}
Given data points $(x_i, y_i)$, $i = 0, 1, \ldots, m$, and a family of functions $\{g(x; a_0, \ldots, a_n)\}$ with $n < m$, the \textbf{approximation problem} consists of finding the parameters $a_0, \ldots, a_n$ that minimize a measure of the discrepancy between $g(x_i; a_0, \ldots, a_n)$ and $y_i$.
\end{definition}

The two principal approaches are:
\begin{itemize}
    \item \textbf{Least squares approximation}: minimize the sum of squared residuals;
    \item \textbf{Chebyshev (minimax) approximation}: minimize the maximum absolute error.
\end{itemize}

%==============================================================================
\section{Discrete Least Squares}
%==============================================================================

\subsection{Problem formulation}

\begin{definition}[Discrete least squares problem]
\label{def:least-squares}
\index{least squares!discrete}
Given $m+1$ data points $(x_i, y_i)$ for $i = 0, \ldots, m$ and a basis of functions $\varphi_0, \varphi_1, \ldots, \varphi_n$ with $n < m$, we seek coefficients $a_0, a_1, \ldots, a_n$ minimizing
\[
    S(a_0, \ldots, a_n) = \sum_{i=0}^{m} \left( y_i - \sum_{j=0}^{n} a_j \varphi_j(x_i) \right)^2.
\]
\end{definition}

\begin{remark}
The condition $n < m$ is essential: the system is \emph{overdetermined}, so we cannot in general satisfy all equations exactly. When $n = m$, we recover the interpolation problem.
\end{remark}

\subsection{Normal equations}

\begin{theorem}[Normal equations]
\label{thm:normal-equations}
\index{normal equations}
The least squares solution satisfies the \textbf{normal equations}:
\[
    \mathbf{A}^\top \mathbf{A} \, \mathbf{a} = \mathbf{A}^\top \mathbf{y},
\]
where $\mathbf{A}$ is the $(m+1) \times (n+1)$ matrix with entries $A_{ij} = \varphi_j(x_i)$, $\mathbf{a} = (a_0, \ldots, a_n)^\top$, and $\mathbf{y} = (y_0, \ldots, y_m)^\top$.
\end{theorem}

\begin{proof}
Define $\mathbf{r} = \mathbf{y} - \mathbf{A}\mathbf{a}$ (the residual vector). Then
\[
    S(\mathbf{a}) = \|\mathbf{r}\|_2^2 = (\mathbf{y} - \mathbf{A}\mathbf{a})^\top (\mathbf{y} - \mathbf{A}\mathbf{a})
    = \mathbf{y}^\top\mathbf{y} - 2\mathbf{a}^\top\mathbf{A}^\top\mathbf{y} + \mathbf{a}^\top\mathbf{A}^\top\mathbf{A}\,\mathbf{a}.
\]
Setting the gradient with respect to $\mathbf{a}$ to zero:
\[
    \nabla_{\mathbf{a}} S = -2\mathbf{A}^\top\mathbf{y} + 2\mathbf{A}^\top\mathbf{A}\,\mathbf{a} = \mathbf{0}.
\]
This gives $\mathbf{A}^\top\mathbf{A}\,\mathbf{a} = \mathbf{A}^\top\mathbf{y}$.

Since $S$ is a convex quadratic form (the Hessian $2\mathbf{A}^\top\mathbf{A}$ is positive semi-definite), any critical point is a global minimum. If the columns of $\mathbf{A}$ are linearly independent, then $\mathbf{A}^\top\mathbf{A}$ is positive definite and the solution is unique.
\end{proof}

\begin{example}[Linear least squares fit]
\label{ex:linear-ls}
Given the data
\[
\begin{array}{c|ccccc}
    x_i & 0 & 1 & 2 & 3 & 4 \\
    \hline
    y_i & 1.0 & 2.1 & 2.9 & 4.2 & 4.8
\end{array}
\]
we fit $g(x) = a_0 + a_1 x$. The matrix $\mathbf{A}$ and the normal equations are:
\[
    \mathbf{A} = \begin{pmatrix} 1 & 0 \\ 1 & 1 \\ 1 & 2 \\ 1 & 3 \\ 1 & 4 \end{pmatrix}, \qquad
    \mathbf{A}^\top\mathbf{A} = \begin{pmatrix} 5 & 10 \\ 10 & 30 \end{pmatrix}, \qquad
    \mathbf{A}^\top\mathbf{y} = \begin{pmatrix} 15.0 \\ 38.8 \end{pmatrix}.
\]
Solving: $a_0 = 1.08$, $a_1 = 0.96$, so $g(x) = 1.08 + 0.96x$.
\end{example}

\begin{warning}
The matrix $\mathbf{A}^\top\mathbf{A}$ can be very ill-conditioned, especially when using monomials $\varphi_j(x) = x^j$ of high degree. This is why orthogonal polynomials or QR factorization are preferred in practice.
\end{warning}

%==============================================================================
\subsection{Polynomial least squares via QR factorization}

\begin{proposition}
\label{prop:qr-ls}
\index{QR factorization!least squares}
If $\mathbf{A} = \mathbf{Q}\mathbf{R}$ is the (thin) QR factorization of $\mathbf{A}$ where $\mathbf{Q} \in \R^{(m+1)\times(n+1)}$ has orthonormal columns and $\mathbf{R} \in \R^{(n+1)\times(n+1)}$ is upper triangular, then the least squares solution is
\[
    \mathbf{R}\,\mathbf{a} = \mathbf{Q}^\top \mathbf{y}.
\]
\end{proposition}

\begin{proof}
From $\mathbf{A} = \mathbf{Q}\mathbf{R}$, we get $\mathbf{A}^\top\mathbf{A} = \mathbf{R}^\top\mathbf{Q}^\top\mathbf{Q}\mathbf{R} = \mathbf{R}^\top\mathbf{R}$ and $\mathbf{A}^\top\mathbf{y} = \mathbf{R}^\top\mathbf{Q}^\top\mathbf{y}$. The normal equations become $\mathbf{R}^\top\mathbf{R}\,\mathbf{a} = \mathbf{R}^\top\mathbf{Q}^\top\mathbf{y}$. Since $\mathbf{R}$ is invertible, we may multiply on the left by $(\mathbf{R}^\top)^{-1}$ to obtain $\mathbf{R}\,\mathbf{a} = \mathbf{Q}^\top\mathbf{y}$.
\end{proof}

%==============================================================================
\section{Continuous Least Squares and Orthogonal Polynomials}
%==============================================================================

\subsection{Continuous least squares}

\begin{definition}[Continuous least squares]
\label{def:continuous-ls}
\index{least squares!continuous}
Given a weight function $w(x) > 0$ on $[a,b]$ and a function $f \in C([a,b])$, we seek a polynomial $p_n \in \mathcal{P}_n$ minimizing
\[
    \|f - p_n\|_w^2 = \int_a^b w(x) \left[f(x) - p_n(x)\right]^2 \, dx.
\]
\end{definition}

The solution satisfies the continuous normal equations:
\[
    \sum_{j=0}^{n} a_j \underbrace{\int_a^b w(x)\,\varphi_j(x)\,\varphi_k(x)\,dx}_{\langle \varphi_j, \varphi_k \rangle_w} = \underbrace{\int_a^b w(x)\,f(x)\,\varphi_k(x)\,dx}_{\langle f, \varphi_k \rangle_w}, \qquad k = 0, \ldots, n.
\]

\subsection{Orthogonal polynomials}

\begin{definition}[Orthogonal polynomials]
\label{def:orthogonal-polynomials}
\index{orthogonal polynomials}
A sequence $\{p_k\}_{k \geq 0}$ of polynomials (with $\deg p_k = k$) is \textbf{orthogonal} with respect to the inner product $\langle \cdot, \cdot \rangle_w$ if
\[
    \langle p_j, p_k \rangle_w = \int_a^b w(x)\,p_j(x)\,p_k(x)\,dx = 0 \quad \text{for } j \neq k.
\]
\end{definition}

\begin{theorem}[Three-term recurrence]
\label{thm:three-term}
\index{three-term recurrence}
Orthogonal polynomials satisfy a recurrence of the form
\[
    p_{k+1}(x) = (\alpha_k x + \beta_k)\,p_k(x) - \gamma_k\,p_{k-1}(x), \qquad k \geq 1,
\]
with $p_0(x) = 1$ and $p_1(x) = \alpha_0 x + \beta_0$.
\end{theorem}

\begin{example}[Classical orthogonal polynomials]
\label{ex:classical-orthogonal}
\begin{center}
\begin{tabular}{lccc}
    \toprule
    \textbf{Name} & \textbf{Interval} & $w(x)$ & \textbf{Notation} \\
    \midrule
    Legendre     & $[-1,1]$          & $1$                   & $P_k(x)$ \\
    Chebyshev    & $[-1,1]$          & $(1-x^2)^{-1/2}$     & $T_k(x)$ \\
    Laguerre     & $[0,\infty)$      & $e^{-x}$             & $L_k(x)$ \\
    Hermite      & $(-\infty,\infty)$& $e^{-x^2}$           & $H_k(x)$ \\
    \bottomrule
\end{tabular}
\end{center}
\end{example}

\begin{proposition}
\label{prop:ls-orthogonal}
If $\{p_0, p_1, \ldots, p_n\}$ is an orthogonal basis, the best approximation is
\[
    f_n(x) = \sum_{k=0}^{n} c_k\, p_k(x), \qquad c_k = \frac{\langle f, p_k \rangle_w}{\langle p_k, p_k \rangle_w}.
\]
The normal equations decouple because the Gram matrix is diagonal.
\end{proposition}

%==============================================================================
\section{Chebyshev Polynomials}
%==============================================================================

\begin{definition}[Chebyshev polynomials]
\label{def:chebyshev}
\index{Chebyshev polynomials}
The Chebyshev polynomials of the first kind are defined by
\[
    T_k(x) = \cos(k \arccos x), \qquad x \in [-1,1], \quad k = 0, 1, 2, \ldots
\]
They satisfy the recurrence $T_0(x)=1$, $T_1(x)=x$, $T_{k+1}(x) = 2x\,T_k(x) - T_{k-1}(x)$.
\end{definition}

\begin{theorem}[Properties of Chebyshev polynomials]
\label{thm:chebyshev-properties}
\begin{enumerate}
    \item $T_k$ has exactly $k$ zeros on $(-1,1)$: $x_j = \cos\!\left(\frac{(2j-1)\pi}{2k}\right)$, $j=1,\ldots,k$.
    \item $|T_k(x)| \leq 1$ for all $x \in [-1,1]$.
    \item $T_k$ attains the values $\pm 1$ at $k+1$ extremal points:
    $x_j^* = \cos\!\left(\frac{j\pi}{k}\right)$, $j = 0, \ldots, k$.
    \item Orthogonality: $\displaystyle\int_{-1}^{1} \frac{T_j(x)\,T_k(x)}{\sqrt{1-x^2}}\,dx = \begin{cases} 0 & j \neq k, \\ \pi & j = k = 0, \\ \pi/2 & j = k \neq 0. \end{cases}$
\end{enumerate}
\end{theorem}

\begin{theorem}[Minimax property]
\label{thm:chebyshev-minimax}
\index{minimax property}
Among all monic polynomials of degree $n$, the polynomial $\widetilde{T}_n(x) = 2^{1-n} T_n(x)$ has the smallest uniform norm on $[-1,1]$:
\[
    \max_{x \in [-1,1]} |\widetilde{T}_n(x)| = 2^{1-n}.
\]
\end{theorem}

%==============================================================================
\section{Best Uniform Approximation (Chebyshev / Minimax)}
%==============================================================================

\begin{definition}[Best uniform approximation]
\label{def:best-uniform}
\index{best uniform approximation}
Given $f \in C([a,b])$, the \textbf{best uniform (minimax) approximation} of degree $n$ is the polynomial $p_n^*$ satisfying
\[
    \|f - p_n^*\|_\infty = \min_{p \in \mathcal{P}_n} \|f - p\|_\infty = \min_{p \in \mathcal{P}_n} \max_{x \in [a,b]} |f(x) - p(x)|.
\]
\end{definition}

\begin{theorem}[Chebyshev equioscillation theorem]
\label{thm:equioscillation}
\index{equioscillation theorem}
A polynomial $p_n^* \in \mathcal{P}_n$ is the best uniform approximation to $f \in C([a,b])$ if and only if the error $e(x) = f(x) - p_n^*(x)$ \textbf{equioscillates} at least $n+2$ times: there exist $n+2$ points
\[
    a \leq x_0 < x_1 < \cdots < x_{n+1} \leq b
\]
such that
\[
    e(x_j) = (-1)^j \, \sigma \, \|e\|_\infty, \qquad j = 0, 1, \ldots, n+1,
\]
where $\sigma = \pm 1$.
\end{theorem}

\begin{remark}
The equioscillation theorem guarantees both existence and uniqueness of the best uniform approximation. The Remez algorithm is a classical iterative method for computing it.
\end{remark}

%==============================================================================
\section{Near-best Approximation via Chebyshev Interpolation}
%==============================================================================

\begin{theorem}[Near-optimality of Chebyshev interpolation]
\label{thm:near-best}
Let $p_n^{\mathrm{Cheb}}$ be the interpolation polynomial at Chebyshev nodes and $p_n^*$ the best uniform approximation. Then
\[
    \|f - p_n^{\mathrm{Cheb}}\|_\infty \leq \left(1 + \Lambda_n^{\mathrm{Cheb}}\right) \|f - p_n^*\|_\infty,
\]
where $\Lambda_n^{\mathrm{Cheb}} = \BigO(\ln n)$ is the Lebesgue constant for Chebyshev nodes.
\end{theorem}

\begin{method}
\textbf{Practical guideline:} For most applications, Chebyshev interpolation provides a near-optimal approximation that is much simpler to compute than the true minimax approximation.
\end{method}

%==============================================================================
\section{Python Implementation}
%==============================================================================

\begin{codeblock}[title=\textbf{Python --- Least squares via normal equations and \texttt{polyfit}]}
\begin{minted}{python}
import numpy as np
import matplotlib.pyplot as plt

# Data points
x = np.array([0.0, 1.0, 2.0, 3.0, 4.0])
y = np.array([1.0, 2.1, 2.9, 4.2, 4.8])

# --- Method 1: Normal equations ---
A = np.column_stack([np.ones_like(x), x])
ATA = A.T @ A
ATy = A.T @ y
a_normal = np.linalg.solve(ATA, ATy)
print("Normal equations: a0 = {:.4f}, a1 = {:.4f}".format(*a_normal))

# --- Method 2: QR factorization ---
Q, R = np.linalg.qr(A)
a_qr = np.linalg.solve(R, Q.T @ y)
print("QR method:        a0 = {:.4f}, a1 = {:.4f}".format(*a_qr))

# --- Method 3: numpy.polyfit ---
coeffs = np.polyfit(x, y, 1)  # returns [a1, a0]
print("polyfit:          a0 = {:.4f}, a1 = {:.4f}".format(coeffs[1], coeffs[0]))

# --- Higher degree fit ---
for deg in [1, 2, 3]:
    c = np.polyfit(x, y, deg)
    p = np.poly1d(c)
    residual = np.sqrt(np.sum((y - p(x))**2))
    print(f"Degree {deg}: residual = {residual:.6f}")

# --- Plot ---
xfine = np.linspace(-0.5, 4.5, 200)
plt.figure(figsize=(8, 5))
plt.plot(x, y, 'ko', markersize=8, label='Data')
for deg, ls in [(1, '-'), (2, '--'), (3, ':')]:
    c = np.polyfit(x, y, deg)
    plt.plot(xfine, np.polyval(c, xfine), ls,
             label=f'Degree {deg}')
plt.xlabel('$x$')
plt.ylabel('$y$')
plt.legend()
plt.title('Least squares polynomial fits')
plt.grid(True, alpha=0.3)
plt.tight_layout()
plt.savefig('least_squares_fit.pdf')
plt.show()
\end{minted}
\end{codeblock}

\begin{codeblock}[title=\textbf{Python --- Chebyshev approximation}]
\begin{minted}{python}
import numpy as np
from numpy.polynomial import chebyshev as C

# Approximate f(x) = exp(x) on [-1, 1]
f = lambda x: np.exp(x)
n = 10

# Chebyshev interpolation nodes
k = np.arange(n + 1)
nodes = np.cos((2*k + 1) * np.pi / (2*(n + 1)))
values = f(nodes)

# Chebyshev coefficients
coeffs = C.chebfit(nodes, values, n)

# Evaluate and compare
xfine = np.linspace(-1, 1, 1000)
approx = C.chebval(xfine, coeffs)
error = np.max(np.abs(f(xfine) - approx))
print(f"Max error (degree {n}): {error:.2e}")
\end{minted}
\end{codeblock}

%==============================================================================
\section{Exercises}
%==============================================================================

\begin{exercise}[$\star$]
\label{ex:ch6-ls-line}
Given the data points $(0,1)$, $(1,3)$, $(2,4)$, $(3,5)$, find the least squares line $y = a_0 + a_1 x$.
\begin{enumerate}
    \item Write the matrix $\mathbf{A}$ and the normal equations $\mathbf{A}^\top\mathbf{A}\,\mathbf{a} = \mathbf{A}^\top\mathbf{y}$.
    \item Solve the system.
    \item Compute the residual $\|\mathbf{y} - \mathbf{A}\mathbf{a}\|_2$.
\end{enumerate}
\end{exercise}

\begin{exercise}[$\star$]
\label{ex:ch6-parabola}
Fit a quadratic $y = a_0 + a_1 x + a_2 x^2$ to the data
\[
\begin{array}{c|cccccc}
    x_i & -2 & -1 & 0 & 1 & 2 & 3 \\
    \hline
    y_i & 5.1 & 2.2 & 0.9 & 1.8 & 5.3 & 10.1
\end{array}
\]
using the normal equations. Compare with \texttt{numpy.polyfit}.
\end{exercise}

\begin{exercise}[$\star\star$]
\label{ex:ch6-conditioning}
Consider the Vandermonde matrix $A_{ij} = x_i^j$ for $m+1$ equispaced points on $[0,1]$.
\begin{enumerate}
    \item For $m = 10$ and $n = 2, 4, 6, 8$, compute $\cond_2(\mathbf{A}^\top\mathbf{A})$.
    \item Compare with $\cond_2(\mathbf{R})$ from the QR factorization of $\mathbf{A}$.
    \item What do you conclude about the numerical stability of normal equations?
\end{enumerate}
\end{exercise}

\begin{exercise}[$\star\star$]
\label{ex:ch6-chebyshev-zeros}
\begin{enumerate}
    \item Verify that $T_4(x) = 8x^4 - 8x^2 + 1$ using the recurrence relation.
    \item Compute its four zeros.
    \item Show that $\|T_4\|_\infty = 1$ on $[-1,1]$ and identify the five extremal points.
\end{enumerate}
\end{exercise}

\begin{exercise}[$\star\star\star$]
\label{ex:ch6-remez}
Implement a simplified Remez algorithm:
\begin{enumerate}
    \item Choose $n+2$ initial reference points (e.g., Chebyshev extrema).
    \item Solve the equioscillation system for $p_n$ and the leveled error $E$.
    \item Find the point of maximum error and exchange it with the worst reference point.
    \item Repeat until convergence.
\end{enumerate}
Test on $f(x) = |x|$ on $[-1,1]$ with $n = 2, 4, 6$.
\end{exercise}

\begin{exercise}[$\star\star\star$]
\label{ex:ch6-weighted-ls}
Consider the weighted least squares problem: minimize
\[
    S_w(\mathbf{a}) = \sum_{i=0}^{m} w_i \left(y_i - \sum_{j=0}^{n} a_j \varphi_j(x_i)\right)^2
\]
where $w_i > 0$ are given weights.
\begin{enumerate}
    \item Derive the weighted normal equations.
    \item Show that they can be written as $\mathbf{A}^\top\mathbf{W}\mathbf{A}\,\mathbf{a} = \mathbf{A}^\top\mathbf{W}\mathbf{y}$ where $\mathbf{W} = \diag(w_0, \ldots, w_m)$.
    \item Implement in Python and test with data where some points should be trusted more.
\end{enumerate}
\end{exercise}

%==============================================================================
\section*{Summary of Key Formulas}
%==============================================================================

\begin{keyformulas}
\textbf{Normal equations:}
\[
    \mathbf{A}^\top \mathbf{A} \, \mathbf{a} = \mathbf{A}^\top \mathbf{y}
\]

\textbf{QR approach:}
\[
    \mathbf{R}\,\mathbf{a} = \mathbf{Q}^\top\mathbf{y} \qquad (\mathbf{A} = \mathbf{Q}\mathbf{R})
\]

\textbf{Continuous least squares with orthogonal polynomials:}
\[
    c_k = \frac{\langle f, p_k \rangle_w}{\langle p_k, p_k \rangle_w}, \qquad f_n(x) = \sum_{k=0}^{n} c_k\, p_k(x)
\]

\textbf{Chebyshev polynomials:}
\[
    T_k(x) = \cos(k\arccos x), \qquad T_{k+1}(x) = 2xT_k(x) - T_{k-1}(x)
\]

\textbf{Equioscillation theorem:} $p_n^*$ is the best uniform approximation $\iff$ the error equioscillates at $\geq n+2$ points.

\textbf{Near-best bound:}
\[
    \|f - p_n^{\mathrm{Cheb}}\|_\infty \leq (1 + \Lambda_n^{\mathrm{Cheb}}) \, \|f - p_n^*\|_\infty, \quad \Lambda_n^{\mathrm{Cheb}} = \BigO(\ln n)
\]
\end{keyformulas}

%!TEX root = ../main.tex
% Chapter 7 — Numerical Differentiation and Integration

\chapter{Numerical Differentiation and Integration}
\label{ch:differentiation-integration}


\section{Introduction}

When an analytical expression for a derivative or integral is unavailable---or the function is known only through data points---we resort to \emph{numerical} methods. This chapter covers finite difference formulas for derivatives and the classical quadrature rules for definite integrals.

%==============================================================================
\section{Numerical Differentiation}
%==============================================================================

\subsection{Finite difference formulas}

\begin{definition}[Finite differences]
\label{def:finite-differences}
\index{finite differences}
Given a function $f$ and a step size $h > 0$:
\begin{itemize}
    \item \textbf{Forward difference:} $\displaystyle f'(x) \approx \frac{f(x+h) - f(x)}{h}$
    \item \textbf{Backward difference:} $\displaystyle f'(x) \approx \frac{f(x) - f(x-h)}{h}$
    \item \textbf{Centered difference:} $\displaystyle f'(x) \approx \frac{f(x+h) - f(x-h)}{2h}$
\end{itemize}
\end{definition}

\begin{theorem}[Truncation errors for finite differences]
\label{thm:fd-errors}
\index{truncation error!finite differences}
If $f$ is sufficiently smooth:
\begin{enumerate}
    \item \textbf{Forward/backward difference:}
    \[
        \frac{f(x+h) - f(x)}{h} = f'(x) + \frac{h}{2}f''(\xi_1), \qquad \xi_1 \in (x, x+h).
    \]
    The error is $\BigO(h)$.

    \item \textbf{Centered difference:}
    \[
        \frac{f(x+h) - f(x-h)}{2h} = f'(x) + \frac{h^2}{6}f'''(\xi_2), \qquad \xi_2 \in (x-h, x+h).
    \]
    The error is $\BigO(h^2)$.

    \item \textbf{Second derivative (centered):}
    \[
        \frac{f(x+h) - 2f(x) + f(x-h)}{h^2} = f''(x) + \frac{h^2}{12}f^{(4)}(\xi_3), \qquad \xi_3 \in (x-h, x+h).
    \]
    The error is $\BigO(h^2)$.
\end{enumerate}
\end{theorem}

\begin{proof}
We prove (2). Taylor-expand $f(x+h)$ and $f(x-h)$:
\begin{align*}
    f(x+h) &= f(x) + hf'(x) + \frac{h^2}{2}f''(x) + \frac{h^3}{6}f'''(x) + \BigO(h^4), \\
    f(x-h) &= f(x) - hf'(x) + \frac{h^2}{2}f''(x) - \frac{h^3}{6}f'''(x) + \BigO(h^4).
\end{align*}
Subtracting:
\[
    f(x+h) - f(x-h) = 2hf'(x) + \frac{h^3}{3}f'''(x) + \BigO(h^5).
\]
Dividing by $2h$: $\frac{f(x+h)-f(x-h)}{2h} = f'(x) + \frac{h^2}{6}f'''(x) + \BigO(h^4)$. By the intermediate value theorem, $f'''(x) = f'''(\xi_2)$ for some $\xi_2$.
\end{proof}

\begin{warning}
\textbf{Rounding vs.\ truncation error trade-off.} For machine epsilon $\eps$, the total error in the centered difference is approximately
\[
    E(h) \approx \frac{h^2}{6}|f'''(x)| + \frac{\eps\,|f(x)|}{h}.
\]
The optimal step size is $h^* \approx \eps^{1/3}$, giving error $\approx \eps^{2/3}$.
\end{warning}

\begin{example}
\label{ex:fd-comparison}
For $f(x) = \sin(x)$ at $x = 1$, the exact derivative is $\cos(1) \approx 0.5403023$. Errors:
\begin{center}
\begin{tabular}{ccc}
    \toprule
    $h$ & Forward $\BigO(h)$ & Centered $\BigO(h^2)$ \\
    \midrule
    $10^{-1}$ & $4.6 \times 10^{-2}$ & $1.4 \times 10^{-3}$ \\
    $10^{-2}$ & $4.6 \times 10^{-3}$ & $1.4 \times 10^{-5}$ \\
    $10^{-4}$ & $4.6 \times 10^{-5}$ & $1.4 \times 10^{-9}$ \\
    $10^{-8}$ & $4.6 \times 10^{-9}$ & $6.1 \times 10^{-11}$ \\
    $10^{-12}$& $5.3 \times 10^{-5}$ & $4.2 \times 10^{-5}$ \\
    \bottomrule
\end{tabular}
\end{center}
Note the loss of accuracy for very small $h$ due to round-off.
\end{example}

\subsection{Higher-order formulas}

\begin{proposition}[Fourth-order approximation of $f'$]
\label{prop:fd4}
\[
    f'(x) = \frac{-f(x+2h) + 8f(x+h) - 8f(x-h) + f(x-2h)}{12h} + \BigO(h^4).
\]
\end{proposition}

%==============================================================================
\section{Numerical Integration (Quadrature)}
%==============================================================================

\subsection{The trapezoidal rule}

\begin{definition}[Simple trapezoidal rule]
\label{def:trap-simple}
\index{trapezoidal rule}
\[
    \int_a^b f(x)\,dx \approx T(f) = \frac{b-a}{2}\bigl[f(a) + f(b)\bigr].
\]
\end{definition}

\begin{theorem}[Trapezoidal rule error]
\label{thm:trap-error}
If $f \in C^2([a,b])$,
\[
    \int_a^b f(x)\,dx - T(f) = -\frac{(b-a)^3}{12}\,f''(\xi), \qquad \xi \in (a,b).
\]
\end{theorem}

\begin{proof}
Let $I = [a,b]$ with $h = b-a$. The interpolation polynomial of degree 1 at $a$ and $b$ is
\[
    p_1(x) = f(a)\frac{b-x}{h} + f(b)\frac{x-a}{h}.
\]
The interpolation error is $f(x) - p_1(x) = \frac{f''(\eta(x))}{2}(x-a)(x-b)$ for some $\eta(x) \in (a,b)$. Integrating:
\[
    \int_a^b f(x)\,dx - \frac{h}{2}[f(a)+f(b)] = \int_a^b \frac{f''(\eta(x))}{2}(x-a)(x-b)\,dx.
\]
Since $(x-a)(x-b) \leq 0$ on $[a,b]$, by the mean value theorem for integrals:
\[
    = \frac{f''(\xi)}{2}\int_a^b (x-a)(x-b)\,dx = \frac{f''(\xi)}{2}\cdot\left(-\frac{h^3}{6}\right) = -\frac{h^3}{12}f''(\xi). \qedhere
\]
\end{proof}

\begin{definition}[Composite trapezoidal rule]
\label{def:trap-composite}
\index{trapezoidal rule!composite}
Dividing $[a,b]$ into $n$ equal sub-intervals of width $h = \frac{b-a}{n}$ with $x_k = a + kh$:
\[
    T_n(f) = \frac{h}{2}\left[f(x_0) + 2\sum_{k=1}^{n-1} f(x_k) + f(x_n)\right].
\]
\end{definition}

\begin{theorem}[Composite trapezoidal error]
\label{thm:trap-composite-error}
If $f \in C^2([a,b])$,
\[
    \int_a^b f(x)\,dx - T_n(f) = -\frac{(b-a)h^2}{12}\,f''(\xi) = \BigO(h^2), \qquad \xi \in (a,b).
\]
\end{theorem}

\subsection{Simpson's rule}

\begin{definition}[Simple Simpson's rule]
\label{def:simpson-simple}
\index{Simpson's rule}
\[
    \int_a^b f(x)\,dx \approx S(f) = \frac{b-a}{6}\left[f(a) + 4f\!\left(\frac{a+b}{2}\right) + f(b)\right].
\]
\end{definition}

\begin{theorem}[Simpson's rule error]
\label{thm:simpson-error}
If $f \in C^4([a,b])$,
\[
    \int_a^b f(x)\,dx - S(f) = -\frac{(b-a)^5}{2880}\,f^{(4)}(\xi), \qquad \xi \in (a,b).
\]
Simpson's rule is exact for polynomials of degree $\leq 3$ (one degree of superconvergence).
\end{theorem}

\begin{definition}[Composite Simpson's rule]
\label{def:simpson-composite}
\index{Simpson's rule!composite}
With $n$ even, $h = \frac{b-a}{n}$, and $x_k = a + kh$:
\[
    S_n(f) = \frac{h}{3}\left[f(x_0) + 4\sum_{\substack{k=1 \\ k\text{ odd}}}^{n-1} f(x_k) + 2\sum_{\substack{k=2 \\ k\text{ even}}}^{n-2} f(x_k) + f(x_n)\right].
\]
\end{definition}

\begin{theorem}[Composite Simpson error]
\label{thm:simpson-composite-error}
If $f \in C^4([a,b])$,
\[
    \int_a^b f(x)\,dx - S_n(f) = -\frac{(b-a)h^4}{180}\,f^{(4)}(\xi) = \BigO(h^4), \qquad \xi \in (a,b).
\]
\end{theorem}

\begin{remark}
Comparison of convergence orders:
\begin{center}
\begin{tabular}{lcc}
    \toprule
    \textbf{Rule} & \textbf{Degree of exactness} & \textbf{Order} \\
    \midrule
    Trapezoidal & 1 & $\BigO(h^2)$ \\
    Simpson     & 3 & $\BigO(h^4)$ \\
    \bottomrule
\end{tabular}
\end{center}
\end{remark}

%==============================================================================
\section{Richardson Extrapolation}
%==============================================================================

\begin{definition}[Richardson extrapolation]
\label{def:richardson}
\index{Richardson extrapolation}
If an approximation $A(h)$ satisfies $A(h) = A_0 + c_1 h^p + \BigO(h^{p+q})$, then
\[
    A^*(h) = \frac{2^p A(h/2) - A(h)}{2^p - 1}
\]
is a higher-order approximation: $A^*(h) = A_0 + \BigO(h^{p+q})$.
\end{definition}

\begin{example}[Romberg integration]
\label{ex:romberg}
\index{Romberg integration}
Applying Richardson extrapolation to the composite trapezoidal rule (where the error has only even powers of $h$):
\begin{itemize}
    \item $T(h)$ and $T(h/2)$ combine to give Simpson's rule: $S = \frac{4T(h/2) - T(h)}{3}$;
    \item Further extrapolation yields Boole's rule and higher-order formulas.
\end{itemize}
The Romberg table:
\[
\begin{array}{cccc}
    T_{1,1} & & & \\
    T_{2,1} & T_{2,2} & & \\
    T_{3,1} & T_{3,2} & T_{3,3} & \\
    T_{4,1} & T_{4,2} & T_{4,3} & T_{4,4}
\end{array}
\]
where $T_{i,1} = T_{2^{i-1}}$ (composite trapezoidal with $2^{i-1}$ sub-intervals) and
\[
    T_{i,j} = \frac{4^{j-1}\,T_{i,j-1} - T_{i-1,j-1}}{4^{j-1} - 1}.
\]
\end{example}

%==============================================================================
\section{Convergence Study: $\int_0^1 e^{-x^2}\,dx$}
%==============================================================================

\begin{example}[Convergence table]
\label{ex:convergence-integration}
We compute $I = \int_0^1 e^{-x^2}\,dx \approx 0.746824132812427$ using the composite trapezoidal and Simpson rules with increasing $n$:

\begin{center}
\begin{tabular}{r|cc|cc}
    \toprule
    $n$ & $|I - T_n|$ & ratio & $|I - S_n|$ & ratio \\
    \midrule
    2   & $2.83 \times 10^{-2}$ & ---  & $5.12 \times 10^{-4}$ & --- \\
    4   & $7.18 \times 10^{-3}$ & 3.94 & $3.27 \times 10^{-5}$ & 15.7 \\
    8   & $1.80 \times 10^{-3}$ & 3.99 & $2.06 \times 10^{-6}$ & 15.9 \\
    16  & $4.51 \times 10^{-4}$ & 4.00 & $1.29 \times 10^{-7}$ & 16.0 \\
    32  & $1.13 \times 10^{-4}$ & 4.00 & $8.06 \times 10^{-9}$ & 16.0 \\
    64  & $2.82 \times 10^{-5}$ & 4.00 & $5.04 \times 10^{-10}$& 16.0 \\
    128 & $7.05 \times 10^{-6}$ & 4.00 & $3.15\times 10^{-11}$ & 16.0 \\
    \bottomrule
\end{tabular}
\end{center}
The ratio $|E_n|/|E_{2n}|$ confirms $\BigO(h^2)$ for trapezoidal (ratio $\to 4$) and $\BigO(h^4)$ for Simpson (ratio $\to 16$).
\end{example}

%==============================================================================
\section{Implementations}
%==============================================================================

\begin{codeblock}[title=\textbf{Python --- Numerical integration}]
\begin{minted}{python}
import numpy as np

def composite_trapezoidal(f, a, b, n):
    """Composite trapezoidal rule with n sub-intervals."""
    x = np.linspace(a, b, n + 1)
    h = (b - a) / n
    return h * (0.5*f(x[0]) + np.sum(f(x[1:-1])) + 0.5*f(x[-1]))

def composite_simpson(f, a, b, n):
    """Composite Simpson's rule (n must be even)."""
    assert n % 2 == 0, "n must be even for Simpson's rule"
    x = np.linspace(a, b, n + 1)
    h = (b - a) / n
    return (h/3) * (f(x[0]) + 4*np.sum(f(x[1::2]))
                    + 2*np.sum(f(x[2:-1:2])) + f(x[-1]))

def romberg(f, a, b, max_order=5):
    """Romberg integration."""
    T = np.zeros((max_order, max_order))
    for i in range(max_order):
        n = 2**i
        T[i, 0] = composite_trapezoidal(f, a, b, n)
        for j in range(1, i + 1):
            T[i, j] = (4**j * T[i, j-1] - T[i-1, j-1]) / (4**j - 1)
    return T

# Convergence study
f = lambda x: np.exp(-x**2)
I_exact = 0.746824132812427

print(f"{'n':>5} {'|I-T_n|':>14} {'ratio':>8} "
      f"{'|I-S_n|':>14} {'ratio':>8}")
print("-" * 55)

err_T_prev, err_S_prev = None, None
for k in range(1, 8):
    n = 2**k
    T = composite_trapezoidal(f, 0, 1, n)
    S = composite_simpson(f, 0, 1, n)
    err_T = abs(I_exact - T)
    err_S = abs(I_exact - S)
    ratio_T = err_T_prev / err_T if err_T_prev else float('nan')
    ratio_S = err_S_prev / err_S if err_S_prev else float('nan')
    print(f"{n:5d} {err_T:14.2e} {ratio_T:8.2f} "
          f"{err_S:14.2e} {ratio_S:8.2f}")
    err_T_prev, err_S_prev = err_T, err_S
\end{minted}
\end{codeblock}

\begin{codeblock}[title=\textbf{Julia --- Numerical integration}]
\begin{minted}{julia}
function composite_trapezoidal(f, a, b, n)
    h = (b - a) / n
    x = range(a, b, length=n+1)
    return h * (f(x[1])/2 + sum(f.(x[2:end-1])) + f(x[end])/2)
end

function composite_simpson(f, a, b, n)
    @assert iseven(n) "n must be even"
    h = (b - a) / n
    x = range(a, b, length=n+1)
    return (h/3) * (f(x[1]) + 4*sum(f.(x[2:2:end-1]))
                    + 2*sum(f.(x[3:2:end-2])) + f(x[end]))
end

# Convergence study
f(x) = exp(-x^2)
I_exact = 0.746824132812427

println("  n     |I-T_n|      ratio    |I-S_n|      ratio")
err_T_prev = err_S_prev = NaN
for k in 1:7
    n = 2^k
    T = composite_trapezoidal(f, 0.0, 1.0, n)
    S = composite_simpson(f, 0.0, 1.0, n)
    err_T = abs(I_exact - T)
    err_S = abs(I_exact - S)
    ratio_T = isnan(err_T_prev) ? NaN : err_T_prev / err_T
    ratio_S = isnan(err_S_prev) ? NaN : err_S_prev / err_S
    @printf("%4d  %12.2e  %7.2f  %12.2e  %7.2f\n",
            n, err_T, ratio_T, err_S, ratio_S)
    err_T_prev, err_S_prev = err_T, err_S
end
\end{minted}
\end{codeblock}

%==============================================================================
\section{Exercises}
%==============================================================================

\begin{exercise}[$\star$]
\label{ex:ch7-fd-order}
Use Taylor expansions to derive the truncation error of the forward difference formula and verify that it is $\BigO(h)$.
\end{exercise}

\begin{exercise}[$\star$]
\label{ex:ch7-trap-exact}
Show that the trapezoidal rule is exact for polynomials of degree $\leq 1$.
\end{exercise}

\begin{exercise}[$\star\star$]
\label{ex:ch7-simpson-exact}
Show that Simpson's rule is exact for polynomials of degree $\leq 3$ (not just $\leq 2$).
\textit{Hint:} verify directly for $f(x) = x^3$ on $[a,b]$.
\end{exercise}

\begin{exercise}[$\star\star$]
\label{ex:ch7-optimal-h}
For the centered difference approximation of $f'(x)$:
\begin{enumerate}
    \item Express the total error (truncation + round-off) as a function of $h$.
    \item Find the optimal $h^*$ that minimizes the total error.
    \item Compute $h^*$ and the corresponding error for double precision ($\eps \approx 10^{-16}$).
\end{enumerate}
\end{exercise}

\begin{exercise}[$\star\star$]
\label{ex:ch7-romberg}
Apply Romberg integration to $\int_0^1 \frac{4}{1+x^2}\,dx = \pi$.
Build the Romberg table up to $T_{4,4}$ and compare with $\pi$.
\end{exercise}

\begin{exercise}[$\star\star\star$]
\label{ex:ch7-adaptive}
Implement an adaptive Simpson rule:
\begin{enumerate}
    \item Compute $S_1 = S(a,b)$ and $S_2 = S(a,m) + S(m,b)$ where $m = (a+b)/2$.
    \item If $|S_2 - S_1| < 15\,\varepsilon$ (with $\varepsilon$ a given tolerance), accept $S_2 + (S_2 - S_1)/15$.
    \item Otherwise, recursively refine each half.
\end{enumerate}
Test on $\int_0^1 \sqrt{x}\,dx$ (which has a singularity in the derivative at $x=0$).
\end{exercise}

%==============================================================================
\section*{Summary of Key Formulas}
%==============================================================================

\begin{keyformulas}
\textbf{Finite differences:}
\begin{align*}
    f'(x) &\approx \frac{f(x+h)-f(x)}{h} & & \BigO(h) \\
    f'(x) &\approx \frac{f(x+h)-f(x-h)}{2h} & & \BigO(h^2) \\
    f''(x) &\approx \frac{f(x+h)-2f(x)+f(x-h)}{h^2} & & \BigO(h^2)
\end{align*}

\textbf{Optimal step (centered difference):} $h^* \approx \eps^{1/3}$, error $\approx \eps^{2/3}$.

\textbf{Composite trapezoidal rule:}
\[
    T_n(f) = \frac{h}{2}\left[f(x_0) + 2\sum_{k=1}^{n-1}f(x_k) + f(x_n)\right], \qquad \text{error } = \BigO(h^2)
\]

\textbf{Composite Simpson's rule} ($n$ even):
\[
    S_n(f) = \frac{h}{3}\left[f(x_0) + 4\!\!\sum_{k\text{ odd}} f(x_k) + 2\!\!\sum_{k\text{ even}} f(x_k) + f(x_n)\right], \qquad \text{error } = \BigO(h^4)
\]

\textbf{Richardson extrapolation:}
\[
    A^*(h) = \frac{2^p A(h/2) - A(h)}{2^p - 1}
\]

\textbf{Romberg recurrence:}
\[
    T_{i,j} = \frac{4^{j-1}\,T_{i,j-1} - T_{i-1,j-1}}{4^{j-1}-1}
\]
\end{keyformulas}

%!TEX root = ../main.tex
% Chapter 8 — Gaussian Quadrature

\chapter{Gaussian Quadrature}
\label{ch:gauss-quadrature}


\section{Introduction}

Newton--Cotes formulas (trapezoidal, Simpson) use \emph{equally spaced} nodes. A natural question arises: \emph{can we do better by choosing the nodes optimally?} Gaussian quadrature answers positively: with $n$ nodes, one can integrate exactly all polynomials of degree up to $2n - 1$.

%==============================================================================
\section{Fundamental Theorem of Gaussian Quadrature}
%==============================================================================

\begin{definition}[Quadrature rule]
\label{def:quadrature}
\index{quadrature rule}
A quadrature rule on $[a,b]$ with weight function $w(x) > 0$ is an approximation
\[
    \int_a^b w(x)\,f(x)\,dx \approx Q_n(f) = \sum_{i=1}^{n} w_i\, f(x_i),
\]
where $x_1, \ldots, x_n$ are the \textbf{nodes} and $w_1, \ldots, w_n$ are the \textbf{weights}.
\end{definition}

\begin{definition}[Degree of exactness]
\label{def:degree-exactness}
\index{degree of exactness}
A quadrature rule has \textbf{degree of exactness} $d$ if it integrates exactly all polynomials of degree $\leq d$ but not all polynomials of degree $d+1$.
\end{definition}

\begin{theorem}[Maximum degree of exactness]
\label{thm:max-exactness}
An $n$-point quadrature rule can have degree of exactness at most $2n-1$. This maximum is achieved if and only if the nodes are the zeros of the degree-$n$ orthogonal polynomial with respect to $w(x)$ on $[a,b]$.
\end{theorem}

\begin{theorem}[Gauss quadrature --- existence and properties]
\label{thm:gauss-quadrature}
\index{Gauss quadrature}
Let $\{p_k\}_{k \geq 0}$ be the sequence of orthogonal polynomials with respect to $\langle f, g \rangle = \int_a^b w(x)\,f(x)\,g(x)\,dx$. Then:
\begin{enumerate}
    \item The $n$ zeros of $p_n$ are real, distinct, and lie in $(a,b)$.
    \item Using these zeros $x_1, \ldots, x_n$ as nodes, there exist unique weights $w_1, \ldots, w_n$ such that $Q_n$ has degree of exactness $2n-1$.
    \item All weights are positive: $w_i > 0$ for $i = 1, \ldots, n$.
\end{enumerate}
\end{theorem}

\begin{proof}[Proof sketch]
\textbf{(1)} Let $x_{i_1}, \ldots, x_{i_k}$ ($k \leq n$) be the zeros of $p_n$ in $(a,b)$ where $p_n$ changes sign. Define $q(x) = (x - x_{i_1})\cdots(x-x_{i_k})$. Then $p_n(x)\,q(x) \geq 0$ on $(a,b)$ (or $\leq 0$), so
\[
    \int_a^b w(x)\,p_n(x)\,q(x)\,dx \neq 0.
\]
But if $k < n$, then $\deg q < n$ and orthogonality of $p_n$ to lower-degree polynomials gives $\int_a^b w(x)\,p_n(x)\,q(x)\,dx = 0$, a contradiction. Hence $k = n$: all $n$ zeros are in $(a,b)$ and $p_n$ changes sign at each.

\textbf{(2)} Define the weights via Lagrange interpolation at the nodes:
\[
    w_i = \int_a^b w(x)\,\ell_i(x)\,dx, \qquad \ell_i(x) = \prod_{\substack{j=1 \\ j\neq i}}^{n} \frac{x - x_j}{x_i - x_j}.
\]
For any polynomial $f$ of degree $\leq 2n-1$, write $f = p_n \cdot q + r$ with $\deg q, \deg r \leq n-1$. Then:
\[
    \int_a^b w(x)\,f(x)\,dx = \underbrace{\int_a^b w(x)\,p_n(x)\,q(x)\,dx}_{=\,0 \text{ (orthogonality)}} + \int_a^b w(x)\,r(x)\,dx.
\]
Since $\deg r \leq n-1$, the interpolation quadrature integrates $r$ exactly, and $f(x_i) = r(x_i)$ since $p_n(x_i) = 0$. Thus $Q_n(f) = Q_n(r) = \int_a^b w(x)\,r(x)\,dx = \int_a^b w(x)\,f(x)\,dx$.

\textbf{(3)} Set $f(x) = \ell_i(x)^2$ (degree $2n-2 \leq 2n-1$). Then $Q_n(f) = w_i > 0$ since $\int_a^b w(x)\,\ell_i(x)^2\,dx > 0$.
\end{proof}

%==============================================================================
\section{Gauss--Legendre Quadrature}
%==============================================================================

\begin{definition}[Gauss--Legendre quadrature]
\label{def:gauss-legendre}
\index{Gauss--Legendre quadrature}
For $w(x) = 1$ on $[-1,1]$, the Gauss--Legendre quadrature uses the zeros of the Legendre polynomial $P_n(x)$:
\[
    \int_{-1}^{1} f(x)\,dx \approx \sum_{i=1}^{n} w_i\, f(x_i).
\]
\end{definition}

\begin{example}[Gauss--Legendre nodes and weights]
\label{ex:gauss-legendre-table}
\begin{center}
\begin{tabular}{c|l|l|c}
    \toprule
    $n$ & Nodes $x_i$ & Weights $w_i$ & Exactness \\
    \midrule
    1 & $0$ & $2$ & $d = 1$ \\[4pt]
    2 & $\pm 1/\sqrt{3}$ & $1$ & $d = 3$ \\[4pt]
    3 & $0, \;\pm\sqrt{3/5}$ & $8/9, \;5/9$ & $d = 5$ \\[4pt]
    4 & $\pm\sqrt{\frac{3-2\sqrt{6/5}}{7}}, \;\pm\sqrt{\frac{3+2\sqrt{6/5}}{7}}$ & $\frac{18+\sqrt{30}}{36}, \;\frac{18-\sqrt{30}}{36}$ & $d = 7$ \\[4pt]
    5 & $0, \;\pm\frac{1}{3}\sqrt{5-2\sqrt{10/7}}, \;\pm\frac{1}{3}\sqrt{5+2\sqrt{10/7}}$ & $\frac{128}{225}, \;\frac{322+13\sqrt{70}}{900}, \;\frac{322-13\sqrt{70}}{900}$ & $d = 9$ \\
    \bottomrule
\end{tabular}
\end{center}
\end{example}

\subsection{Change of interval}

\begin{proposition}[Interval transformation]
\label{prop:interval-change}
\index{interval change}
To apply Gauss--Legendre quadrature on a general interval $[a,b]$, use the linear transformation $x = \frac{b-a}{2}\,t + \frac{a+b}{2}$:
\[
    \int_a^b f(x)\,dx = \frac{b-a}{2}\int_{-1}^{1} f\!\left(\frac{b-a}{2}\,t + \frac{a+b}{2}\right)dt \approx \frac{b-a}{2}\sum_{i=1}^{n} w_i\, f\!\left(\frac{b-a}{2}\,t_i + \frac{a+b}{2}\right).
\]
\end{proposition}

%==============================================================================
\section{Error Formula}
%==============================================================================

\begin{theorem}[Gauss quadrature error]
\label{thm:gauss-error}
\index{Gauss quadrature!error}
If $f \in C^{2n}([a,b])$, the error of the $n$-point Gauss quadrature is
\[
    E_n(f) = \int_a^b w(x)\,f(x)\,dx - \sum_{i=1}^{n} w_i\,f(x_i) = \frac{f^{(2n)}(\xi)}{(2n)!}\int_a^b w(x)\,\left[\prod_{i=1}^{n}(x-x_i)\right]^2 dx
\]
for some $\xi \in (a,b)$.
\end{theorem}

\begin{corollary}[Gauss--Legendre error]
For $n$-point Gauss--Legendre quadrature on $[-1,1]$:
\[
    E_n(f) = \frac{2^{2n+1}(n!)^4}{(2n+1)[(2n)!]^3}\,f^{(2n)}(\xi), \qquad \xi \in (-1,1).
\]
\end{corollary}

%==============================================================================
\section{Other Gaussian Quadrature Rules}
%==============================================================================

\subsection{Gauss--Laguerre}

\begin{definition}[Gauss--Laguerre quadrature]
\label{def:gauss-laguerre}
\index{Gauss--Laguerre quadrature}
For integrals of the form $\int_0^\infty e^{-x}\,f(x)\,dx$, the nodes are the zeros of the Laguerre polynomial $L_n(x)$:
\[
    \int_0^\infty e^{-x}\,f(x)\,dx \approx \sum_{i=1}^{n} w_i\,f(x_i).
\]
\end{definition}

\begin{example}[Gauss--Laguerre nodes and weights ($n=3$)]
\begin{center}
\begin{tabular}{ccc}
    \toprule
    $i$ & $x_i$ & $w_i$ \\
    \midrule
    1 & $0.415775$ & $0.711093$ \\
    2 & $2.294280$ & $0.278518$ \\
    3 & $6.289945$ & $0.010389$ \\
    \bottomrule
\end{tabular}
\end{center}
\end{example}

\subsection{Gauss--Hermite}

\begin{definition}[Gauss--Hermite quadrature]
\label{def:gauss-hermite}
\index{Gauss--Hermite quadrature}
For integrals of the form $\int_{-\infty}^{\infty} e^{-x^2}\,f(x)\,dx$, the nodes are the zeros of the Hermite polynomial $H_n(x)$:
\[
    \int_{-\infty}^{\infty} e^{-x^2}\,f(x)\,dx \approx \sum_{i=1}^{n} w_i\,f(x_i).
\]
\end{definition}

\begin{example}[Gauss--Hermite nodes and weights ($n=3$)]
\begin{center}
\begin{tabular}{ccc}
    \toprule
    $i$ & $x_i$ & $w_i$ \\
    \midrule
    1 & $-\sqrt{3/2}$ & $\sqrt{\pi}/6$ \\
    2 & $0$           & $2\sqrt{\pi}/3$ \\
    3 & $\sqrt{3/2}$  & $\sqrt{\pi}/6$ \\
    \bottomrule
\end{tabular}
\end{center}
\end{example}

\begin{remark}
Summary of Gaussian quadrature variants:
\begin{center}
\begin{tabular}{llcl}
    \toprule
    \textbf{Type} & \textbf{Interval} & $w(x)$ & \textbf{Polynomials} \\
    \midrule
    Gauss--Legendre  & $[-1,1]$            & $1$          & Legendre $P_n$ \\
    Gauss--Chebyshev & $[-1,1]$            & $(1-x^2)^{-1/2}$ & Chebyshev $T_n$ \\
    Gauss--Laguerre  & $[0,\infty)$        & $e^{-x}$     & Laguerre $L_n$ \\
    Gauss--Hermite   & $(-\infty,\infty)$  & $e^{-x^2}$   & Hermite $H_n$ \\
    \bottomrule
\end{tabular}
\end{center}
\end{remark}

%==============================================================================
\section{Composite Gauss Quadrature}
%==============================================================================

\begin{method}
To improve accuracy without using a large number of nodes per element, subdivide $[a,b]$ into panels and apply Gauss quadrature on each panel. With $m$ panels of equal width $h = (b-a)/m$ and $n$-point Gauss quadrature on each:
\[
    \int_a^b f(x)\,dx \approx \sum_{k=1}^{m} \frac{h}{2} \sum_{i=1}^{n} w_i\, f\!\left(\frac{h}{2}\,t_i + \frac{x_{k-1}+x_k}{2}\right).
\]
The overall error is $\BigO(h^{2n})$.
\end{method}

%==============================================================================
\section{Implementations}
%==============================================================================

\begin{codeblock}[title=\textbf{Python --- Gauss--Legendre quadrature}]
\begin{minted}{python}
import numpy as np
from numpy.polynomial.legendre import leggauss

def gauss_legendre(f, a, b, n):
    """n-point Gauss-Legendre quadrature on [a, b]."""
    nodes, weights = leggauss(n)
    # Transform from [-1, 1] to [a, b]
    t = 0.5 * (b - a) * nodes + 0.5 * (a + b)
    return 0.5 * (b - a) * np.sum(weights * f(t))

# Test: integral of exp(-x^2) from 0 to 1
f = lambda x: np.exp(-x**2)
I_exact = 0.746824132812427

print(f"{'n':>4} {'Approximation':>18} {'Error':>14}")
print("-" * 40)
for n in [2, 3, 4, 5, 8, 10, 15, 20]:
    approx = gauss_legendre(f, 0, 1, n)
    err = abs(I_exact - approx)
    print(f"{n:4d} {approx:18.15f} {err:14.2e}")

# Composite Gauss-Legendre
def composite_gauss(f, a, b, n_nodes, n_panels):
    """Composite Gauss-Legendre quadrature."""
    h = (b - a) / n_panels
    result = 0.0
    for k in range(n_panels):
        ak = a + k * h
        bk = ak + h
        result += gauss_legendre(f, ak, bk, n_nodes)
    return result

# Compare convergence
print("\nComposite Gauss-Legendre (3-point per panel):")
for m in [1, 2, 4, 8, 16]:
    approx = composite_gauss(f, 0, 1, 3, m)
    err = abs(I_exact - approx)
    print(f"  panels = {m:3d}, error = {err:.2e}")
\end{minted}
\end{codeblock}

\begin{codeblock}[title=\textbf{Julia --- Gauss--Legendre quadrature with FastGaussQuadrature}]
\begin{minted}{julia}
using FastGaussQuadrature

function gauss_legendre(f, a, b, n)
    nodes, weights = gausslegendre(n)
    # Transform from [-1, 1] to [a, b]
    t = @. 0.5 * (b - a) * nodes + 0.5 * (a + b)
    return 0.5 * (b - a) * sum(weights .* f.(t))
end

# Test
f(x) = exp(-x^2)
I_exact = 0.746824132812427

println("   n      Approximation          Error")
println("-" ^ 45)
for n in [2, 3, 4, 5, 8, 10, 15, 20]
    approx = gauss_legendre(f, 0.0, 1.0, n)
    err = abs(I_exact - approx)
    @printf("%4d  %18.15f  %14.2e\n", n, approx, err)
end

# Gauss-Laguerre
println("\nGauss-Laguerre for integral of x^2 * exp(-x):")
for n in [3, 5, 10]
    nodes, weights = gausslaguerre(n)
    approx = sum(weights .* nodes.^2)  # exact = 2
    @printf("  n = %2d, approx = %.15f, error = %.2e\n",
            n, approx, abs(2.0 - approx))
end

# Gauss-Hermite
println("\nGauss-Hermite for integral of x^4 * exp(-x^2):")
for n in [3, 5, 10]
    nodes, weights = gausshermite(n)
    approx = sum(weights .* nodes.^4)  # exact = 3*sqrt(pi)/4
    exact = 3*sqrt(pi)/4
    @printf("  n = %2d, approx = %.15f, error = %.2e\n",
            n, approx, abs(exact - approx))
end
\end{minted}
\end{codeblock}

%==============================================================================
\section{Exercises}
%==============================================================================

\begin{exercise}[$\star$]
\label{ex:ch8-gl2}
Verify that the $2$-point Gauss--Legendre rule ($x_{1,2} = \pm 1/\sqrt{3}$, $w_{1,2} = 1$) is exact for $f(x) = x^3$ on $[-1,1]$.
\end{exercise}

\begin{exercise}[$\star$]
\label{ex:ch8-interval}
Use 3-point Gauss--Legendre quadrature to approximate $\int_0^2 e^x\,dx$. Compare with the exact value $e^2 - 1$.
\end{exercise}

\begin{exercise}[$\star\star$]
\label{ex:ch8-derive-gl2}
Derive the $2$-point Gauss--Legendre rule from scratch by requiring exactness for $f(x) = 1, x, x^2, x^3$. Solve the resulting nonlinear system for $x_1, x_2, w_1, w_2$.
\end{exercise}

\begin{exercise}[$\star\star$]
\label{ex:ch8-composite}
\begin{enumerate}
    \item Implement composite Gauss--Legendre quadrature (3-point per panel).
    \item Compute $\int_0^1 e^{-x^2}\,dx$ with $m = 1, 2, 4, 8, 16$ panels.
    \item Verify numerically that the error is $\BigO(h^6)$ where $h = 1/m$.
\end{enumerate}
\end{exercise}

\begin{exercise}[$\star\star$]
\label{ex:ch8-laguerre}
Use Gauss--Laguerre quadrature to compute $\Gamma(s) = \int_0^\infty x^{s-1}\,e^{-x}\,dx$ for $s = 1.5, 2.5, 3.5$. Compare with the exact values $\Gamma(1.5) = \frac{\sqrt{\pi}}{2}$, $\Gamma(2.5) = \frac{3\sqrt{\pi}}{4}$, $\Gamma(3.5) = \frac{15\sqrt{\pi}}{8}$.
\end{exercise}

\begin{exercise}[$\star\star\star$]
\label{ex:ch8-eigenvalue-method}
The Gauss quadrature nodes and weights can be computed from the eigenvalue problem of the tridiagonal Jacobi matrix.
\begin{enumerate}
    \item For Gauss--Legendre, construct the symmetric tridiagonal matrix with $\alpha_k = 0$ and $\beta_k = k/\sqrt{4k^2-1}$.
    \item Compute its eigenvalues (the nodes) and eigenvectors (to obtain the weights: $w_i = 2\,v_{1,i}^2$ where $v_{1,i}$ is the first component of the $i$-th eigenvector).
    \item Compare with \texttt{numpy.polynomial.legendre.leggauss} for $n = 5, 10, 20$.
\end{enumerate}
\end{exercise}

%==============================================================================
\section*{Summary of Key Formulas}
%==============================================================================

\begin{keyformulas}
\textbf{Gauss quadrature:}
\[
    \int_a^b w(x)\,f(x)\,dx \approx \sum_{i=1}^{n} w_i\,f(x_i)
\]
Nodes = zeros of the degree-$n$ orthogonal polynomial; degree of exactness $= 2n-1$.

\textbf{Gauss--Legendre on $[-1,1]$} ($w(x) = 1$, Legendre polynomials):
\[
    \int_{-1}^{1} f(x)\,dx \approx \sum_{i=1}^{n} w_i\,f(x_i), \qquad w_i > 0, \quad \sum_{i=1}^n w_i = 2
\]

\textbf{Change of interval:}
\[
    \int_a^b f(x)\,dx = \frac{b-a}{2}\int_{-1}^{1} f\!\left(\tfrac{b-a}{2}\,t + \tfrac{a+b}{2}\right)dt
\]

\textbf{Error formula:}
\[
    E_n(f) = \frac{f^{(2n)}(\xi)}{(2n)!}\int_a^b w(x)\left[\prod_{i=1}^{n}(x-x_i)\right]^2 dx
\]

\textbf{Summary table:}
\begin{center}
\begin{tabular}{lccc}
    Gauss--Legendre & $[-1,1]$ & $w=1$ & Legendre $P_n$ \\
    Gauss--Laguerre & $[0,\infty)$ & $e^{-x}$ & Laguerre $L_n$ \\
    Gauss--Hermite  & $(-\infty,\infty)$ & $e^{-x^2}$ & Hermite $H_n$ \\
\end{tabular}
\end{center}
\end{keyformulas}

%!TEX root = ../main.tex
% Chapter 9 — Numerical Solution of ODEs

\chapter{Numerical Solution of Ordinary Differential Equations}
\label{ch:ode}


\section{Introduction}

Many phenomena in physics, biology, and engineering are modeled by ordinary differential equations (ODEs) that cannot be solved analytically. This chapter introduces the fundamental numerical methods for initial value problems.

%==============================================================================
\section{The Cauchy Problem}
%==============================================================================

\begin{definition}[Initial value problem / Cauchy problem]
\label{def:cauchy}
\index{Cauchy problem}
\index{initial value problem}
Given $f: [t_0, T] \times \R^d \to \R^d$, find $y: [t_0, T] \to \R^d$ such that
\[
    \begin{cases}
        y'(t) = f(t, y(t)), & t \in [t_0, T], \\
        y(t_0) = y_0.
    \end{cases}
\]
\end{definition}

\begin{theorem}[Picard--Lindel\"of]
\label{thm:picard}
\index{Picard--Lindel\"of theorem}
If $f$ is continuous in $t$ and Lipschitz continuous in $y$ with Lipschitz constant $L$:
\[
    \|f(t,y_1) - f(t,y_2)\| \leq L\,\|y_1 - y_2\| \quad \forall\, t \in [t_0, T],\; y_1, y_2 \in \R^d,
\]
then the Cauchy problem has a unique solution on $[t_0, T]$.
\end{theorem}

We discretize $[t_0, T]$ with a uniform grid: $t_n = t_0 + nh$, $h = (T-t_0)/N$, and seek approximations $y_n \approx y(t_n)$.

%==============================================================================
\section{Euler's Methods}
%==============================================================================

\subsection{Explicit (forward) Euler method}

\begin{definition}[Explicit Euler method]
\label{def:euler-explicit}
\index{Euler method!explicit}
\[
    y_{n+1} = y_n + h\,f(t_n, y_n), \qquad n = 0, 1, \ldots, N-1.
\]
\end{definition}

\begin{algorithm}[H]
\caption{Explicit Euler method}
\label{alg:euler-explicit}
\KwInput{$f$, $t_0$, $T$, $y_0$, $N$ (number of steps)}
\KwOutput{Arrays $t[0:N]$, $y[0:N]$}
$h \gets (T - t_0) / N$\;
$t[0] \gets t_0$\;
$y[0] \gets y_0$\;
\For{$n = 0, 1, \ldots, N-1$}{
    $y[n+1] \gets y[n] + h \cdot f(t[n], y[n])$\;
    $t[n+1] \gets t[n] + h$\;
}
\Return{$t$, $y$}\;
\end{algorithm}

\begin{theorem}[Global error of explicit Euler]
\label{thm:euler-error}
\index{Euler method!error}
If $f$ is Lipschitz in $y$ with constant $L$ and $y \in C^2([t_0,T])$, then
\[
    \max_{0 \leq n \leq N} \|y(t_n) - y_n\| \leq \frac{hM}{2L}\left(e^{L(T-t_0)} - 1\right),
\]
where $M = \max_{t \in [t_0,T]} \|y''(t)\|$. The method is of \textbf{order 1}: the global error is $\BigO(h)$.
\end{theorem}

\begin{proof}
\textbf{Local truncation error.} Taylor-expanding:
\[
    y(t_{n+1}) = y(t_n) + h\,y'(t_n) + \frac{h^2}{2}\,y''(\xi_n) = y(t_n) + h\,f(t_n, y(t_n)) + \frac{h^2}{2}\,y''(\xi_n).
\]
The local truncation error is $\tau_n = \frac{h}{2}\,y''(\xi_n)$, so $|\tau_n| \leq \frac{hM}{2}$.

\textbf{Global error.} Let $e_n = y(t_n) - y_n$. Then:
\begin{align*}
    e_{n+1} &= y(t_{n+1}) - y_{n+1} \\
    &= \bigl[y(t_n) + h\,f(t_n, y(t_n)) + h\tau_n\bigr] - \bigl[y_n + h\,f(t_n, y_n)\bigr] \\
    &= e_n + h\bigl[f(t_n, y(t_n)) - f(t_n, y_n)\bigr] + h\tau_n.
\end{align*}
Using the Lipschitz condition:
\[
    \|e_{n+1}\| \leq (1 + hL)\|e_n\| + h\,\frac{hM}{2}.
\]
Setting $\delta = \frac{hM}{2}$ and iterating with $e_0 = 0$:
\[
    \|e_n\| \leq \frac{\delta}{hL}\bigl[(1+hL)^n - 1\bigr] \leq \frac{hM}{2L}\bigl[e^{nhL} - 1\bigr] \leq \frac{hM}{2L}\bigl[e^{L(T-t_0)} - 1\bigr],
\]
using the inequality $(1+hL)^n \leq e^{nhL}$.
\end{proof}

\subsection{Implicit (backward) Euler method}

\begin{definition}[Implicit Euler method]
\label{def:euler-implicit}
\index{Euler method!implicit}
\[
    y_{n+1} = y_n + h\,f(t_{n+1}, y_{n+1}), \qquad n = 0, 1, \ldots, N-1.
\]
This is an implicit equation for $y_{n+1}$ and generally requires solving a nonlinear system at each step (e.g., by Newton's method).
\end{definition}

\begin{definition}[A-stability]
\label{def:A-stability}
\index{A-stability}
A numerical method is \textbf{A-stable} if its stability region contains the entire left half of the complex plane: $\{z \in \C : \operatorname{Re}(z) \leq 0\} \subset \mathcal{S}$.
\end{definition}

\begin{theorem}
\label{thm:implicit-euler-Astable}
The implicit Euler method is A-stable.
\end{theorem}

\begin{proof}
For the test equation $y' = \lambda y$ ($\lambda \in \C$), the implicit Euler gives:
\[
    y_{n+1} = y_n + h\lambda\,y_{n+1} \implies y_{n+1} = \frac{y_n}{1 - h\lambda}.
\]
The amplification factor is $R(z) = \frac{1}{1-z}$ where $z = h\lambda$. Stability requires $|R(z)| \leq 1$, i.e., $|1-z| \geq 1$. For $z = x + iy$ with $x \leq 0$: $|1-z|^2 = (1-x)^2 + y^2 \geq 1$ since $1-x \geq 1$. Thus $|R(z)| \leq 1$ for all $\operatorname{Re}(z) \leq 0$.
\end{proof}

%==============================================================================
\section{Runge--Kutta Methods}
%==============================================================================

\begin{definition}[General $s$-stage Runge--Kutta method]
\label{def:runge-kutta}
\index{Runge--Kutta method}
An $s$-stage Runge--Kutta method is defined by
\begin{align*}
    k_i &= f\!\left(t_n + c_i h,\; y_n + h\sum_{j=1}^{s} a_{ij}\,k_j\right), \qquad i = 1, \ldots, s, \\
    y_{n+1} &= y_n + h\sum_{i=1}^{s} b_i\,k_i.
\end{align*}
The method is \textbf{explicit} if $a_{ij} = 0$ for $j \geq i$.
\end{definition}

\begin{definition}[Butcher tableau]
\label{def:butcher}
\index{Butcher tableau}
A Runge--Kutta method is compactly represented by its \textbf{Butcher tableau}:
\[
\renewcommand{\arraystretch}{1.2}
\begin{array}{c|c}
    \mathbf{c} & \mathbf{A} \\
    \hline
    & \mathbf{b}^\top
\end{array}
\qquad = \qquad
\begin{array}{c|cccc}
    c_1 & a_{11} & a_{12} & \cdots & a_{1s} \\
    c_2 & a_{21} & a_{22} & \cdots & a_{2s} \\
    \vdots & \vdots & \vdots & \ddots & \vdots \\
    c_s & a_{s1} & a_{s2} & \cdots & a_{ss} \\
    \hline
    & b_1 & b_2 & \cdots & b_s
\end{array}
\]
with the consistency condition $c_i = \sum_{j=1}^{s} a_{ij}$.
\end{definition}

\subsection{Classical fourth-order Runge--Kutta (RK4)}

\begin{definition}[RK4]
\label{def:rk4}
\index{Runge--Kutta method!RK4}
The classical RK4 method:
\begin{align*}
    k_1 &= f(t_n, y_n), \\
    k_2 &= f(t_n + h/2, \; y_n + (h/2)\,k_1), \\
    k_3 &= f(t_n + h/2, \; y_n + (h/2)\,k_2), \\
    k_4 &= f(t_n + h, \; y_n + h\,k_3), \\
    y_{n+1} &= y_n + \frac{h}{6}(k_1 + 2k_2 + 2k_3 + k_4).
\end{align*}
Butcher tableau:
\[
\begin{array}{c|cccc}
    0   &     &     &     &   \\
    1/2 & 1/2 &     &     &   \\
    1/2 & 0   & 1/2 &     &   \\
    1   & 0   & 0   & 1   &   \\
    \hline
        & 1/6 & 1/3 & 1/3 & 1/6
\end{array}
\]
\end{definition}

\begin{theorem}[Order of RK4]
\label{thm:rk4-order}
The classical RK4 method has \textbf{order 4}: the global error is $\BigO(h^4)$ and the local truncation error is $\BigO(h^5)$.
\end{theorem}

\begin{remark}
For explicit Runge--Kutta methods:
\begin{center}
\begin{tabular}{lcc}
    \toprule
    \textbf{Method} & \textbf{Stages} & \textbf{Order} \\
    \midrule
    Euler       & 1 & 1 \\
    Heun (RK2)  & 2 & 2 \\
    Kutta (RK3) & 3 & 3 \\
    Classic RK4 & 4 & 4 \\
    \bottomrule
\end{tabular}
\end{center}
For $s \geq 5$ stages, the maximum achievable order is less than $s$ (Butcher barrier).
\end{remark}

%==============================================================================
\section{Stability Analysis}
%==============================================================================

\begin{definition}[Test equation]
\label{def:test-equation}
\index{test equation}
The Dahlquist test equation is $y' = \lambda y$, $y(0) = 1$, with $\lambda \in \C$. Any one-step method applied to this equation produces
\[
    y_{n+1} = R(z)\,y_n, \qquad z = h\lambda,
\]
where $R(z)$ is the \textbf{amplification factor} (or stability function).
\end{definition}

\begin{definition}[Stability region]
\label{def:stability-region}
\index{stability region}
The \textbf{stability region} of a method is
\[
    \mathcal{S} = \{z \in \C : |R(z)| \leq 1\}.
\]
For the exact solution $|e^{z}| \leq 1$ when $\operatorname{Re}(z) \leq 0$, so a method is A-stable if $\mathcal{S} \supseteq \{z : \operatorname{Re}(z) \leq 0\}$.
\end{definition}

\begin{proposition}[Amplification factors]
\label{prop:amplification}
\index{amplification factor}
\begin{center}
\begin{tabular}{ll}
    \toprule
    \textbf{Method} & $R(z)$ \\
    \midrule
    Explicit Euler  & $1 + z$ \\
    Implicit Euler  & $\dfrac{1}{1-z}$ \\
    Trapezoidal (Crank--Nicolson) & $\dfrac{1+z/2}{1-z/2}$ \\
    RK4             & $1 + z + \dfrac{z^2}{2} + \dfrac{z^3}{6} + \dfrac{z^4}{24}$ \\
    \bottomrule
\end{tabular}
\end{center}
\end{proposition}

\begin{remark}
The explicit Euler stability region is the disk $|1+z| \leq 1$, centered at $(-1,0)$ with radius $1$. For the explicit Euler applied to $y' = \lambda y$ with $\lambda < 0$, stability requires $0 < h < 2/|\lambda|$.
\end{remark}

\subsection{Stability region plots}

\begin{figure}[htbp]
\centering
\begin{tikzpicture}[scale=1.3]
    % Axes
    \draw[->] (-4,0) -- (2,0) node[right] {$\operatorname{Re}(z)$};
    \draw[->] (0,-3.5) -- (0,3.5) node[above] {$\operatorname{Im}(z)$};

    % Explicit Euler: |1+z| <= 1, circle centered at (-1,0) radius 1
    \draw[blue, thick, fill=blue!10, fill opacity=0.3] (-1,0) circle (1);
    \node[blue] at (-1,-1.4) {\small Explicit Euler};

    % RK4 stability region (approximate boundary)
    \draw[red, thick, fill=red!10, fill opacity=0.2]
        plot[smooth cycle, tension=0.7] coordinates {
            (-2.78,0) (-2.6,0.6) (-2.2,1.2) (-1.5,1.8)
            (-0.8,2.4) (-0.1,2.8) (0.2,2.5) (0.3,1.5)
            (0.1,0.5) (0.1,-0.5) (0.3,-1.5) (0.2,-2.5)
            (-0.1,-2.8) (-0.8,-2.4) (-1.5,-1.8) (-2.2,-1.2)
            (-2.6,-0.6)
        };
    \node[red] at (-1.5,2.5) {\small RK4};

    % Grid
    \foreach \x in {-3,-2,-1,1}
        \draw (\x,0.05) -- (\x,-0.05) node[below] {\tiny $\x$};
    \foreach \y in {-3,-2,-1,1,2,3}
        \draw (0.05,\y) -- (-0.05,\y) node[left] {\tiny $\y i$};
\end{tikzpicture}
\caption{Stability regions of explicit Euler (blue) and RK4 (red) in the complex $z = h\lambda$ plane.}
\label{fig:stability-regions}
\end{figure}

%==============================================================================
\section{Convergence Study: $y' = -y$}
%==============================================================================

\begin{example}[Convergence table]
\label{ex:ode-convergence}
We solve $y' = -y$, $y(0) = 1$ on $[0,1]$ (exact: $y(1) = e^{-1}$):

\begin{center}
\begin{tabular}{r|cc|cc|cc}
    \toprule
    $N$ & $|e_N^{\text{Euler}}|$ & order & $|e_N^{\text{Heun}}|$ & order & $|e_N^{\text{RK4}}|$ & order \\
    \midrule
    10   & $1.7\times10^{-2}$ & ---  & $4.1\times10^{-4}$ & ---  & $1.5\times10^{-7}$ & --- \\
    20   & $8.5\times10^{-3}$ & 1.0  & $1.0\times10^{-4}$ & 2.0  & $9.3\times10^{-9}$ & 4.0 \\
    40   & $4.2\times10^{-3}$ & 1.0  & $2.6\times10^{-5}$ & 2.0  & $5.8\times10^{-10}$& 4.0 \\
    80   & $2.1\times10^{-3}$ & 1.0  & $6.4\times10^{-6}$ & 2.0  & $3.6\times10^{-11}$& 4.0 \\
    160  & $1.1\times10^{-3}$ & 1.0  & $1.6\times10^{-6}$ & 2.0  & $2.3\times10^{-12}$& 4.0 \\
    \bottomrule
\end{tabular}
\end{center}
\end{example}

%==============================================================================
\section{Implementations}
%==============================================================================

\begin{codeblock}[title=\textbf{Python --- ODE solvers}]
\begin{minted}{python}
import numpy as np
import matplotlib.pyplot as plt

def euler_explicit(f, t0, T, y0, N):
    """Explicit Euler method."""
    h = (T - t0) / N
    t = np.linspace(t0, T, N + 1)
    y = np.zeros((N + 1, len(np.atleast_1d(y0))))
    y[0] = y0
    for n in range(N):
        y[n+1] = y[n] + h * f(t[n], y[n])
    return t, y

def euler_implicit(f, jac, t0, T, y0, N, tol=1e-12):
    """Implicit Euler method (Newton iteration)."""
    h = (T - t0) / N
    t = np.linspace(t0, T, N + 1)
    d = len(np.atleast_1d(y0))
    y = np.zeros((N + 1, d))
    y[0] = y0
    for n in range(N):
        # Newton iteration: G(u) = u - y[n] - h*f(t[n+1], u) = 0
        u = y[n].copy()  # initial guess
        for _ in range(50):
            G = u - y[n] - h * f(t[n+1], u)
            dG = np.eye(d) - h * jac(t[n+1], u)
            delta = np.linalg.solve(dG, -G)
            u += delta
            if np.linalg.norm(delta) < tol:
                break
        y[n+1] = u
    return t, y

def rk4(f, t0, T, y0, N):
    """Classical RK4 method."""
    h = (T - t0) / N
    t = np.linspace(t0, T, N + 1)
    y = np.zeros((N + 1, len(np.atleast_1d(y0))))
    y[0] = y0
    for n in range(N):
        k1 = f(t[n], y[n])
        k2 = f(t[n] + h/2, y[n] + h/2 * k1)
        k3 = f(t[n] + h/2, y[n] + h/2 * k2)
        k4 = f(t[n] + h, y[n] + h * k3)
        y[n+1] = y[n] + (h/6) * (k1 + 2*k2 + 2*k3 + k4)
    return t, y

# Test: y' = -y, y(0) = 1
f = lambda t, y: -y
y0 = np.array([1.0])
y_exact = lambda t: np.exp(-t)

print(f"{'N':>5} {'Euler':>14} {'order':>6} "
      f"{'RK4':>14} {'order':>6}")
print("-" * 50)
err_e_prev = err_r_prev = None
for N in [10, 20, 40, 80, 160, 320]:
    t_e, y_e = euler_explicit(f, 0, 1, y0, N)
    t_r, y_r = rk4(f, 0, 1, y0, N)
    err_e = abs(y_e[-1, 0] - y_exact(1.0))
    err_r = abs(y_r[-1, 0] - y_exact(1.0))
    ord_e = np.log2(err_e_prev/err_e) if err_e_prev else float('nan')
    ord_r = np.log2(err_r_prev/err_r) if err_r_prev else float('nan')
    print(f"{N:5d} {err_e:14.2e} {ord_e:6.2f} "
          f"{err_r:14.2e} {ord_r:6.2f}")
    err_e_prev, err_r_prev = err_e, err_r
\end{minted}
\end{codeblock}

\begin{codeblock}[title=\textbf{Julia --- ODE solvers}]
\begin{minted}{julia}
function euler_explicit(f, t0, T, y0, N)
    h = (T - t0) / N
    t = range(t0, T, length=N+1)
    d = length(y0)
    y = zeros(d, N+1)
    y[:, 1] = y0
    for n in 1:N
        y[:, n+1] = y[:, n] + h * f(t[n], y[:, n])
    end
    return collect(t), y
end

function rk4(f, t0, T, y0, N)
    h = (T - t0) / N
    t = range(t0, T, length=N+1)
    d = length(y0)
    y = zeros(d, N+1)
    y[:, 1] = y0
    for n in 1:N
        k1 = f(t[n], y[:, n])
        k2 = f(t[n] + h/2, y[:, n] + h/2 * k1)
        k3 = f(t[n] + h/2, y[:, n] + h/2 * k2)
        k4 = f(t[n] + h, y[:, n] + h * k3)
        y[:, n+1] = y[:, n] + (h/6) * (k1 + 2k2 + 2k3 + k4)
    end
    return collect(t), y
end

# Test: y' = -y
f(t, y) = -y
y0 = [1.0]

for N in [10, 20, 40, 80, 160]
    t_e, y_e = euler_explicit(f, 0.0, 1.0, y0, N)
    t_r, y_r = rk4(f, 0.0, 1.0, y0, N)
    err_e = abs(y_e[1, end] - exp(-1.0))
    err_r = abs(y_r[1, end] - exp(-1.0))
    @printf("N = %4d: Euler = %.2e, RK4 = %.2e\n",
            N, err_e, err_r)
end
\end{minted}
\end{codeblock}

%==============================================================================
\section{Exercises}
%==============================================================================

\begin{exercise}[$\star$]
\label{ex:ch9-euler-hand}
Apply the explicit Euler method with $h = 0.25$ to solve $y' = -2y$, $y(0)=1$ on $[0,1]$. Compare with the exact solution $y(t) = e^{-2t}$.
\end{exercise}

\begin{exercise}[$\star$]
\label{ex:ch9-stability-euler}
For the test equation $y' = \lambda y$ with $\lambda = -10$:
\begin{enumerate}
    \item Determine the maximum step size $h$ for which the explicit Euler method is stable.
    \item Show that the implicit Euler method is stable for any $h > 0$.
\end{enumerate}
\end{exercise}

\begin{exercise}[$\star\star$]
\label{ex:ch9-rk2-butcher}
The midpoint method has Butcher tableau:
\[
\begin{array}{c|cc}
    0   &     &   \\
    1/2 & 1/2 &   \\
    \hline
        & 0   & 1
\end{array}
\]
\begin{enumerate}
    \item Write out the method explicitly.
    \item Show it has order 2 by verifying the order conditions.
    \item Determine its stability function $R(z)$ and plot the stability region.
\end{enumerate}
\end{exercise}

\begin{exercise}[$\star\star$]
\label{ex:ch9-convergence-table}
Reproduce the convergence table of Example~\ref{ex:ode-convergence}. Verify numerically the convergence orders of explicit Euler, Heun's method, and RK4.
\end{exercise}

\begin{exercise}[$\star\star$]
\label{ex:ch9-stiff}
Consider the stiff system
\[
    y' = -1000(y - \sin t) + \cos t, \qquad y(0) = 0.
\]
\begin{enumerate}
    \item Solve with explicit Euler for $h = 0.001$ and $h = 0.003$. What happens?
    \item Solve with implicit Euler for $h = 0.01$ and $h = 0.1$.
    \item Explain the results using stability analysis.
\end{enumerate}
\end{exercise}

\begin{exercise}[$\star\star\star$ --- Project: Lorenz system]
\label{ex:ch9-lorenz}
\index{Lorenz system}
The Lorenz system is
\[
    \begin{cases}
        \dot{x} = \sigma(y - x), \\
        \dot{y} = x(\rho - z) - y, \\
        \dot{z} = xy - \beta z,
    \end{cases}
\]
with the classical parameters $\sigma = 10$, $\rho = 28$, $\beta = 8/3$.
\begin{enumerate}
    \item Implement the system and solve with RK4 for $t \in [0, 50]$ with $h = 0.01$.
    \item Plot the trajectory in 3D ($x$--$y$--$z$ phase space).
    \item Investigate sensitivity to initial conditions: solve with $(x_0, y_0, z_0) = (1, 1, 1)$ and $(1.001, 1, 1)$. Plot $|x_1(t) - x_2(t)|$ and observe exponential divergence.
    \item Estimate the largest Lyapunov exponent.
\end{enumerate}
\end{exercise}

\begin{exercise}[$\star\star\star$]
\label{ex:ch9-adaptive}
Implement an adaptive RK4-5 method (Dormand--Prince):
\begin{enumerate}
    \item Use embedded RK pairs of order 4 and 5 to estimate the local error.
    \item Adjust the step size: $h_{\text{new}} = h \cdot \min\!\left(2,\; \max\!\left(0.5,\; 0.9\left(\frac{\text{tol}}{|e|}\right)^{1/5}\right)\right)$.
    \item Test on $y' = -y^2$, $y(0) = 1$ (exact: $y = 1/(1+t)$) with tolerance $10^{-8}$.
\end{enumerate}
\end{exercise}

%==============================================================================
\section*{Summary of Key Formulas}
%==============================================================================

\begin{keyformulas}
\textbf{Explicit Euler:}
\[
    y_{n+1} = y_n + h\,f(t_n, y_n), \qquad \text{order } 1, \quad R(z) = 1 + z
\]

\textbf{Implicit Euler:}
\[
    y_{n+1} = y_n + h\,f(t_{n+1}, y_{n+1}), \qquad \text{order } 1, \quad R(z) = \frac{1}{1-z}, \quad \text{A-stable}
\]

\textbf{RK4:}
\[
    y_{n+1} = y_n + \frac{h}{6}(k_1 + 2k_2 + 2k_3 + k_4), \qquad \text{order } 4
\]

\textbf{Global error bound (Euler):}
\[
    \max_n \|y(t_n) - y_n\| \leq \frac{hM}{2L}(e^{L(T-t_0)}-1)
\]

\textbf{Stability condition (explicit Euler):} $z = h\lambda \in \{z : |1+z| \leq 1\}$

\textbf{A-stability:} stability region $\supseteq \{z \in \C : \operatorname{Re}(z) \leq 0\}$

\textbf{Butcher tableau:}
\[
\begin{array}{c|c}
    \mathbf{c} & \mathbf{A} \\
    \hline
    & \mathbf{b}^\top
\end{array}
\]
\end{keyformulas}

%!TEX root = ../main.tex
% Chapter 10 — Eigenvalues and Eigenvectors

\chapter{Numerical Computation of Eigenvalues and Eigenvectors}
\label{ch:eigenvalues}


\section{Introduction}

The eigenvalue problem---finding $\lambda$ and $\mathbf{v} \neq \mathbf{0}$ such that $\mathbf{A}\mathbf{v} = \lambda\mathbf{v}$---arises throughout science and engineering: vibration analysis, stability of dynamical systems, principal component analysis, Google's PageRank, and quantum mechanics, among many others. Since the eigenvalues are roots of the characteristic polynomial $\det(\mathbf{A} - \lambda\mathbf{I}) = 0$, and there is no closed-form formula for roots of polynomials of degree $\geq 5$ (Abel--Ruffini theorem), all general eigenvalue algorithms are necessarily \emph{iterative}.

%==============================================================================
\section{Spectral Properties}
%==============================================================================

\begin{definition}[Spectrum]
\label{def:spectrum}
\index{spectrum}
The \textbf{spectrum} of $\mathbf{A} \in \R^{n \times n}$ (or $\C^{n \times n}$) is the set of its eigenvalues:
\[
    \sigma(\mathbf{A}) = \{\lambda \in \C : \det(\mathbf{A} - \lambda\mathbf{I}) = 0\}.
\]
The \textbf{spectral radius} is $\rho(\mathbf{A}) = \max_{\lambda \in \sigma(\mathbf{A})} |\lambda|$.
\end{definition}

\begin{proposition}[Basic spectral properties]
\label{prop:spectral}
Let $\mathbf{A} \in \C^{n \times n}$ with eigenvalues $\lambda_1, \ldots, \lambda_n$ (counted with multiplicity).
\begin{enumerate}
    \item $\displaystyle\sum_{i=1}^n \lambda_i = \operatorname{tr}(\mathbf{A})$, \qquad $\displaystyle\prod_{i=1}^n \lambda_i = \det(\mathbf{A})$.
    \item If $\mathbf{A}$ is real and $\lambda \in \sigma(\mathbf{A})$, then $\bar{\lambda} \in \sigma(\mathbf{A})$.
    \item If $\mathbf{A}$ is symmetric (Hermitian), all eigenvalues are real.
    \item If $\mathbf{A}$ is symmetric positive definite, all eigenvalues are positive.
    \item $\sigma(\mathbf{A}^{-1}) = \{1/\lambda : \lambda \in \sigma(\mathbf{A})\}$ when $\mathbf{A}$ is invertible.
    \item $\sigma(\mathbf{A} - \sigma\mathbf{I}) = \{\lambda - \sigma : \lambda \in \sigma(\mathbf{A})\}$ (spectral shift).
\end{enumerate}
\end{proposition}

%==============================================================================
\section{Gerschgorin's Theorem}
%==============================================================================

\begin{theorem}[Gerschgorin circle theorem]
\label{thm:gerschgorin}
\index{Gerschgorin theorem}
Every eigenvalue of $\mathbf{A} \in \C^{n \times n}$ lies in at least one of the \textbf{Gerschgorin discs}:
\[
    D_i = \left\{ z \in \C : |z - a_{ii}| \leq R_i \right\}, \qquad R_i = \sum_{\substack{j=1 \\ j \neq i}}^{n} |a_{ij}|.
\]
That is, $\sigma(\mathbf{A}) \subset \bigcup_{i=1}^{n} D_i$.

Moreover, if a union of $k$ discs is disjoint from the remaining $n-k$ discs, it contains exactly $k$ eigenvalues (counted with multiplicity).
\end{theorem}

\begin{proof}
Let $\lambda$ be an eigenvalue with eigenvector $\mathbf{v}$, and let $i$ be the index of the component of largest modulus: $|v_i| = \|\mathbf{v}\|_\infty$. From $\mathbf{A}\mathbf{v} = \lambda\mathbf{v}$, the $i$-th equation gives:
\[
    \sum_{j=1}^{n} a_{ij}\,v_j = \lambda\,v_i \implies (\lambda - a_{ii})\,v_i = \sum_{\substack{j=1 \\ j \neq i}}^{n} a_{ij}\,v_j.
\]
Taking absolute values:
\[
    |\lambda - a_{ii}|\,|v_i| \leq \sum_{\substack{j=1 \\ j \neq i}}^{n} |a_{ij}|\,|v_j| \leq \left(\sum_{\substack{j=1 \\ j \neq i}}^{n} |a_{ij}|\right)|v_i| = R_i\,|v_i|.
\]
Since $|v_i| > 0$, we get $|\lambda - a_{ii}| \leq R_i$, i.e., $\lambda \in D_i$.
\end{proof}

\begin{example}[Gerschgorin discs]
\label{ex:gerschgorin}
For the matrix
\[
    \mathbf{A} = \begin{pmatrix} 4 & 1 & 0 \\ 0.5 & 3 & 0.5 \\ 0 & 1 & 6 \end{pmatrix},
\]
the Gerschgorin discs are:
\begin{itemize}
    \item $D_1$: center $4$, radius $1$, i.e., $[3, 5]$;
    \item $D_2$: center $3$, radius $1$, i.e., $[2, 4]$;
    \item $D_3$: center $6$, radius $1$, i.e., $[5, 7]$.
\end{itemize}
The actual eigenvalues are approximately $2.59$, $4.25$, $6.16$, all within the union of discs. Since $D_3$ is disjoint from $D_1 \cup D_2$, it contains exactly one eigenvalue.
\end{example}

\subsection{Gerschgorin discs visualization}

\begin{figure}[htbp]
\centering
\begin{tikzpicture}[scale=1.0]
    % Axes
    \draw[->] (0.5,0) -- (8.5,0) node[right] {$\operatorname{Re}$};
    \draw[->] (1,-2) -- (1,2) node[above] {$\operatorname{Im}$};

    % Disc D1: center 4, radius 1
    \draw[blue, thick, fill=blue!10, fill opacity=0.4] (4,0) circle (1);
    \node[blue, above] at (4,1.1) {$D_1$};
    \fill[blue] (4,0) circle (2pt);

    % Disc D2: center 3, radius 1
    \draw[red, thick, fill=red!10, fill opacity=0.4] (3,0) circle (1);
    \node[red, below] at (3,-1.2) {$D_2$};
    \fill[red] (3,0) circle (2pt);

    % Disc D3: center 6, radius 1
    \draw[green!50!black, thick, fill=green!10, fill opacity=0.4] (6,0) circle (1);
    \node[green!50!black, above] at (6,1.1) {$D_3$};
    \fill[green!50!black] (6,0) circle (2pt);

    % Eigenvalues (approximate)
    \fill[black] (2.59,0) circle (3pt) node[below] {\small $\lambda_1$};
    \fill[black] (4.25,0) circle (3pt) node[below right] {\small $\lambda_2$};
    \fill[black] (6.16,0) circle (3pt) node[below] {\small $\lambda_3$};

    % Axis labels
    \foreach \x in {2,3,4,5,6,7}
        \draw (\x,0.07) -- (\x,-0.07) node[below] {\small $\x$};
\end{tikzpicture}
\caption{Gerschgorin discs for the matrix $\mathbf{A}$. Black dots indicate the actual eigenvalues.}
\label{fig:gerschgorin}
\end{figure}

%==============================================================================
\section{Power Iteration}
%==============================================================================

\begin{definition}[Power iteration]
\label{def:power-iteration}
\index{power iteration}
The \textbf{power iteration} (or power method) computes the dominant eigenvalue (largest in modulus) and its eigenvector.
\end{definition}

\begin{algorithm}[H]
\caption{Power iteration}
\label{alg:power-iteration}
\KwInput{$\mathbf{A} \in \R^{n \times n}$, initial vector $\mathbf{v}^{(0)}$, tolerance $\varepsilon$, max iterations $k_{\max}$}
\KwOutput{Dominant eigenvalue $\lambda$, eigenvector $\mathbf{v}$}
$\mathbf{v} \gets \mathbf{v}^{(0)} / \|\mathbf{v}^{(0)}\|_2$\;
\For{$k = 1, 2, \ldots, k_{\max}$}{
    $\mathbf{w} \gets \mathbf{A}\,\mathbf{v}$\;
    $\lambda \gets \mathbf{v}^\top \mathbf{w}$ \quad \tcp{Rayleigh quotient}
    $\mathbf{v} \gets \mathbf{w} / \|\mathbf{w}\|_2$\;
    \If{$\|\mathbf{A}\mathbf{v} - \lambda\mathbf{v}\|_2 < \varepsilon$}{
        \Return{$\lambda$, $\mathbf{v}$}\;
    }
}
\Return{$\lambda$, $\mathbf{v}$}\;
\end{algorithm}

\begin{theorem}[Convergence of power iteration]
\label{thm:power-convergence}
\index{power iteration!convergence}
Let $\mathbf{A} \in \R^{n \times n}$ have eigenvalues $|\lambda_1| > |\lambda_2| \geq \cdots \geq |\lambda_n|$ with corresponding eigenvectors $\mathbf{v}_1, \ldots, \mathbf{v}_n$ forming a basis. If $\mathbf{v}^{(0)}$ has a nonzero component in the direction of $\mathbf{v}_1$, then:
\begin{enumerate}
    \item The Rayleigh quotient converges to $\lambda_1$:
    \[
        |\lambda^{(k)} - \lambda_1| = \BigO\!\left(\left|\frac{\lambda_2}{\lambda_1}\right|^{2k}\right).
    \]
    \item The iterates $\mathbf{v}^{(k)}$ converge to $\pm\mathbf{v}_1$ (up to normalization):
    \[
        \|\mathbf{v}^{(k)} - (\pm\mathbf{v}_1)\| = \BigO\!\left(\left|\frac{\lambda_2}{\lambda_1}\right|^k\right).
    \]
\end{enumerate}
The rate of convergence is $r = |\lambda_2/\lambda_1| < 1$.
\end{theorem}

\begin{proof}
Write the initial vector as $\mathbf{v}^{(0)} = \sum_{i=1}^{n} c_i\,\mathbf{v}_i$ with $c_1 \neq 0$. After $k$ iterations (ignoring normalization):
\[
    \mathbf{A}^k\,\mathbf{v}^{(0)} = \sum_{i=1}^{n} c_i\,\lambda_i^k\,\mathbf{v}_i = \lambda_1^k\left[c_1\,\mathbf{v}_1 + \sum_{i=2}^{n} c_i\left(\frac{\lambda_i}{\lambda_1}\right)^k \mathbf{v}_i\right].
\]
Since $|\lambda_i/\lambda_1| < 1$ for $i \geq 2$, the sum converges to $c_1\lambda_1^k\,\mathbf{v}_1$, so the normalized iterate converges to $\pm\mathbf{v}_1$.

For the Rayleigh quotient, $\lambda^{(k)} = (\mathbf{v}^{(k)})^\top\mathbf{A}\,\mathbf{v}^{(k)}$. Since the eigenvector error is $\BigO(r^k)$, the eigenvalue error is $\BigO(r^{2k})$ (the Rayleigh quotient has a quadratic stationary property for symmetric matrices).
\end{proof}

\begin{remark}
The convergence can be very slow when $|\lambda_2/\lambda_1|$ is close to 1. This motivates the use of shifted methods and more sophisticated algorithms.
\end{remark}

%==============================================================================
\section{Inverse Iteration with Shift}
%==============================================================================

\begin{definition}[Inverse iteration with shift]
\label{def:inverse-iteration}
\index{inverse iteration}
To find the eigenvalue closest to a given shift $\sigma$, we apply power iteration to $(\mathbf{A} - \sigma\mathbf{I})^{-1}$, whose dominant eigenvalue is $1/(\lambda_j - \sigma)$ where $\lambda_j$ is the eigenvalue of $\mathbf{A}$ closest to $\sigma$.
\end{definition}

\begin{algorithm}[H]
\caption{Inverse iteration with shift}
\label{alg:inverse-iteration}
\KwInput{$\mathbf{A} \in \R^{n\times n}$, shift $\sigma$, initial $\mathbf{v}^{(0)}$, tolerance $\varepsilon$, $k_{\max}$}
\KwOutput{Eigenvalue $\lambda$ closest to $\sigma$, eigenvector $\mathbf{v}$}
$\mathbf{v} \gets \mathbf{v}^{(0)} / \|\mathbf{v}^{(0)}\|_2$\;
Compute LU factorization of $\mathbf{A} - \sigma\mathbf{I}$\;
\For{$k = 1, 2, \ldots, k_{\max}$}{
    Solve $(\mathbf{A} - \sigma\mathbf{I})\,\mathbf{w} = \mathbf{v}$ \quad \tcp{using LU factors}
    $\lambda \gets \mathbf{v}^\top\mathbf{A}\,\mathbf{v}$ \quad \tcp{Rayleigh quotient}
    $\mathbf{v} \gets \mathbf{w} / \|\mathbf{w}\|_2$\;
    \If{$\|\mathbf{A}\mathbf{v} - \lambda\mathbf{v}\|_2 < \varepsilon$}{
        \Return{$\lambda$, $\mathbf{v}$}\;
    }
}
\Return{$\lambda$, $\mathbf{v}$}\;
\end{algorithm}

\begin{proposition}
Inverse iteration with shift $\sigma$ converges at rate $\left|\frac{\lambda_j - \sigma}{\lambda_k - \sigma}\right|$, where $\lambda_j$ is the nearest eigenvalue to $\sigma$ and $\lambda_k$ is the second nearest. With a good estimate of $\sigma$, convergence can be very fast.
\end{proposition}

\begin{remark}[Rayleigh quotient iteration]
\index{Rayleigh quotient iteration}
If the shift $\sigma$ is updated at each step to the current Rayleigh quotient, we obtain the \textbf{Rayleigh quotient iteration}, which converges cubically for symmetric matrices.
\end{remark}

%==============================================================================
\section{The QR Algorithm}
%==============================================================================

\subsection{Basic QR algorithm}

\begin{definition}[QR algorithm]
\label{def:qr-algorithm}
\index{QR algorithm}
Starting from $\mathbf{A}_0 = \mathbf{A}$, the basic QR algorithm iterates:
\[
    \mathbf{A}_k = \mathbf{Q}_k\,\mathbf{R}_k \qquad (\text{QR factorization}), \qquad
    \mathbf{A}_{k+1} = \mathbf{R}_k\,\mathbf{Q}_k.
\]
\end{definition}

\begin{proposition}[Properties of the QR iteration]
\label{prop:qr-properties}
\begin{enumerate}
    \item All matrices $\mathbf{A}_k$ are similar: $\mathbf{A}_{k+1} = \mathbf{Q}_k^\top\mathbf{A}_k\,\mathbf{Q}_k$.
    \item If $\mathbf{A}$ has eigenvalues $|\lambda_1| > |\lambda_2| > \cdots > |\lambda_n| > 0$, then $\mathbf{A}_k$ converges to an upper triangular matrix with the eigenvalues on the diagonal.
    \item The convergence rate of the sub-diagonal entry $(i+1,i)$ is $\BigO(|\lambda_{i+1}/\lambda_i|^k)$.
\end{enumerate}
\end{proposition}

\subsection{QR algorithm with Hessenberg reduction}

\begin{definition}[Hessenberg form]
\label{def:hessenberg}
\index{Hessenberg form}
A matrix is in \textbf{upper Hessenberg form} if $a_{ij} = 0$ for $i > j+1$:
\[
    \mathbf{H} = \begin{pmatrix}
        h_{11} & h_{12} & h_{13} & \cdots & h_{1n} \\
        h_{21} & h_{22} & h_{23} & \cdots & h_{2n} \\
        0      & h_{32} & h_{33} & \cdots & h_{3n} \\
        \vdots & \ddots & \ddots & \ddots & \vdots \\
        0      & \cdots & 0      & h_{n,n-1} & h_{nn}
    \end{pmatrix}.
\]
\end{definition}

\begin{theorem}[Hessenberg reduction]
\label{thm:hessenberg}
Any matrix $\mathbf{A} \in \R^{n \times n}$ can be reduced to upper Hessenberg form by an orthogonal similarity transformation: $\mathbf{H} = \mathbf{Q}^\top\mathbf{A}\,\mathbf{Q}$, using Householder reflections, in $\BigO(n^3)$ operations.
\end{theorem}

\begin{method}
The practical QR algorithm proceeds in two phases:
\begin{enumerate}
    \item \textbf{Reduction:} Transform $\mathbf{A}$ to Hessenberg form $\mathbf{H}$ (cost: $\BigO(n^3)$, done once).
    \item \textbf{Iteration:} Apply the QR iteration to $\mathbf{H}$, which preserves Hessenberg form. Each QR step on a Hessenberg matrix costs only $\BigO(n^2)$.
\end{enumerate}
For symmetric matrices, Hessenberg form is tridiagonal, and each step costs $\BigO(n)$.
\end{method}

\subsection{QR algorithm with shifts}

\begin{definition}[Shifted QR iteration]
\label{def:shifted-qr}
\index{QR algorithm!shift}
\[
    \mathbf{A}_k - \sigma_k \mathbf{I} = \mathbf{Q}_k\,\mathbf{R}_k, \qquad
    \mathbf{A}_{k+1} = \mathbf{R}_k\,\mathbf{Q}_k + \sigma_k \mathbf{I}.
\]
The shift $\sigma_k$ is chosen to accelerate convergence.
\end{definition}

\begin{definition}[Wilkinson shift]
\label{def:wilkinson-shift}
\index{Wilkinson shift}
The \textbf{Wilkinson shift} is the eigenvalue of the bottom-right $2 \times 2$ submatrix
\[
    \begin{pmatrix} a_{n-1,n-1} & a_{n-1,n} \\ a_{n,n-1} & a_{nn} \end{pmatrix}
\]
that is closer to $a_{nn}$. Explicitly:
\[
    \sigma = a_{nn} + \delta - \operatorname{sign}(\delta)\sqrt{\delta^2 + a_{n,n-1}^2}, \qquad \delta = \frac{a_{n-1,n-1} - a_{nn}}{2}.
\]
\end{definition}

\begin{theorem}
With the Wilkinson shift, the QR algorithm converges cubically for symmetric matrices and at least quadratically in general.
\end{theorem}

\begin{remark}
The complete modern QR algorithm (as implemented in LAPACK) combines:
\begin{itemize}
    \item Hessenberg reduction;
    \item Implicit double shifts (Francis QR step) to handle complex conjugate eigenvalue pairs in real arithmetic;
    \item Deflation when a sub-diagonal element becomes negligible;
    \item Exceptional shifts for difficult cases.
\end{itemize}
The total cost is $\BigO(n^3)$ for all eigenvalues.
\end{remark}

%==============================================================================
\section{Singular Value Decomposition (SVD)}
%==============================================================================

\begin{theorem}[Singular Value Decomposition]
\label{thm:svd}
\index{singular value decomposition}
\index{SVD|see{singular value decomposition}}
Every matrix $\mathbf{A} \in \R^{m \times n}$ (with $m \geq n$) admits a factorization
\[
    \mathbf{A} = \mathbf{U}\,\boldsymbol{\Sigma}\,\mathbf{V}^\top,
\]
where:
\begin{itemize}
    \item $\mathbf{U} \in \R^{m \times m}$ is orthogonal (columns: left singular vectors),
    \item $\mathbf{V} \in \R^{n \times n}$ is orthogonal (columns: right singular vectors),
    \item $\boldsymbol{\Sigma} \in \R^{m \times n}$ is ``diagonal'' with $\sigma_1 \geq \sigma_2 \geq \cdots \geq \sigma_n \geq 0$ (the \textbf{singular values}).
\end{itemize}
\end{theorem}

\begin{proposition}[Properties of the SVD]
\label{prop:svd-properties}
\begin{enumerate}
    \item $\rang(\mathbf{A}) = $ number of nonzero singular values.
    \item $\|\mathbf{A}\|_2 = \sigma_1$, \quad $\|\mathbf{A}\|_F = \sqrt{\sigma_1^2 + \cdots + \sigma_n^2}$.
    \item $\cond_2(\mathbf{A}) = \sigma_1 / \sigma_n$ (when $\mathbf{A}$ is square and invertible).
    \item The singular values of $\mathbf{A}$ are the square roots of the eigenvalues of $\mathbf{A}^\top\mathbf{A}$.
    \item $\mathbf{A} = \sum_{i=1}^{r} \sigma_i\,\mathbf{u}_i\,\mathbf{v}_i^\top$ (outer product form), where $r = \rang(\mathbf{A})$.
\end{enumerate}
\end{proposition}

\begin{theorem}[Eckart--Young theorem]
\label{thm:eckart-young}
\index{Eckart--Young theorem}
The best rank-$k$ approximation to $\mathbf{A}$ (in both the $2$-norm and the Frobenius norm) is
\[
    \mathbf{A}_k = \sum_{i=1}^{k} \sigma_i\,\mathbf{u}_i\,\mathbf{v}_i^\top = \mathbf{U}_k\,\boldsymbol{\Sigma}_k\,\mathbf{V}_k^\top,
\]
where $\mathbf{U}_k$, $\boldsymbol{\Sigma}_k$, $\mathbf{V}_k$ are the truncated SVD. The approximation errors are:
\[
    \|\mathbf{A} - \mathbf{A}_k\|_2 = \sigma_{k+1}, \qquad
    \|\mathbf{A} - \mathbf{A}_k\|_F = \sqrt{\sigma_{k+1}^2 + \cdots + \sigma_n^2}.
\]
\end{theorem}

%==============================================================================
\section{Convergence Study}
%==============================================================================

\begin{example}[Power iteration convergence]
\label{ex:power-convergence}
For the matrix
\[
    \mathbf{A} = \begin{pmatrix} 4 & 1 & 0 \\ 1 & 3 & 1 \\ 0 & 1 & 2 \end{pmatrix},
\]
the eigenvalues are $\lambda_1 \approx 4.732$, $\lambda_2 \approx 3.000$, $\lambda_3 \approx 1.268$. The ratio $r = |\lambda_2/\lambda_1| \approx 0.634$.

\begin{center}
\begin{tabular}{r|cc|cc}
    \toprule
    $k$ & $\lambda^{(k)}$ & $|\lambda^{(k)} - \lambda_1|$ & $|\lambda_2/\lambda_1|^k$ & $|\lambda_2/\lambda_1|^{2k}$ \\
    \midrule
    1  & $4.5000$ & $2.32 \times 10^{-1}$ & $6.34 \times 10^{-1}$ & $4.02 \times 10^{-1}$ \\
    2  & $4.7143$ & $1.77 \times 10^{-2}$ & $4.02 \times 10^{-1}$ & $1.61 \times 10^{-1}$ \\
    3  & $4.7302$ & $1.76 \times 10^{-3}$ & $2.55 \times 10^{-1}$ & $6.48 \times 10^{-2}$ \\
    5  & $4.7321$ & $1.38 \times 10^{-5}$ & $1.02 \times 10^{-1}$ & $1.04 \times 10^{-2}$ \\
    10 & $4.7321$ & $3.32 \times 10^{-11}$ & $1.05 \times 10^{-2}$ & $1.09 \times 10^{-4}$ \\
    \bottomrule
\end{tabular}
\end{center}
The eigenvalue error decreases at rate $\BigO(r^{2k})$ (Rayleigh quotient).
\end{example}

%==============================================================================
\section{Implementations}
%==============================================================================

\begin{codeblock}[title=\textbf{Python --- Eigenvalue algorithms}]
\begin{minted}{python}
import numpy as np
import matplotlib.pyplot as plt

def power_iteration(A, v0=None, tol=1e-12, maxiter=1000):
    """Power iteration for the dominant eigenvalue."""
    n = A.shape[0]
    v = v0 if v0 is not None else np.random.randn(n)
    v = v / np.linalg.norm(v)
    history = []
    for k in range(maxiter):
        w = A @ v
        lam = v @ w  # Rayleigh quotient
        v = w / np.linalg.norm(w)
        history.append(lam)
        if np.linalg.norm(A @ v - lam * v) < tol:
            break
    return lam, v, history

def inverse_iteration(A, sigma, v0=None, tol=1e-12, maxiter=1000):
    """Inverse iteration with shift sigma."""
    n = A.shape[0]
    v = v0 if v0 is not None else np.random.randn(n)
    v = v / np.linalg.norm(v)
    B = A - sigma * np.eye(n)
    lu_piv = None
    try:
        from scipy.linalg import lu_factor, lu_solve
        lu_piv = lu_factor(B)
    except ImportError:
        pass
    history = []
    for k in range(maxiter):
        if lu_piv is not None:
            w = lu_solve(lu_piv, v)
        else:
            w = np.linalg.solve(B, v)
        lam = v @ (A @ v)  # Rayleigh quotient
        v = w / np.linalg.norm(w)
        history.append(lam)
        if np.linalg.norm(A @ v - lam * v) < tol:
            break
    return lam, v, history

def qr_algorithm(A, maxiter=200, tol=1e-12):
    """Basic QR algorithm (without shifts)."""
    n = A.shape[0]
    Ak = A.copy().astype(float)
    for k in range(maxiter):
        Q, R = np.linalg.qr(Ak)
        Ak = R @ Q
        # Check convergence: sub-diagonal elements
        off = np.sum(np.abs(np.tril(Ak, -1)))
        if off < tol:
            break
    return np.diag(Ak), Ak

# Example
A = np.array([[4, 1, 0],
              [1, 3, 1],
              [0, 1, 2]], dtype=float)

print("=== Power iteration ===")
lam, v, hist = power_iteration(A)
print(f"Dominant eigenvalue: {lam:.10f}")
print(f"Eigenvector: {v}")

print("\n=== Inverse iteration (sigma=1.0) ===")
lam2, v2, hist2 = inverse_iteration(A, sigma=1.0)
print(f"Nearest eigenvalue to 1.0: {lam2:.10f}")

print("\n=== QR algorithm ===")
eigenvalues, _ = qr_algorithm(A)
print(f"All eigenvalues: {np.sort(eigenvalues)[::-1]}")

print("\n=== NumPy reference ===")
eigvals = np.linalg.eigvalsh(A)
print(f"numpy.linalg.eigvalsh: {np.sort(eigvals)[::-1]}")

# SVD example
print("\n=== SVD ===")
M = np.array([[3, 2, 2],
              [2, 3, -2]], dtype=float)
U, s, Vt = np.linalg.svd(M)
print(f"Singular values: {s}")
print(f"Rank: {np.sum(s > 1e-10)}")
print(f"||M||_2 = {s[0]:.6f}")
print(f"cond_2(M) = {s[0]/s[-1]:.6f}")

# Low-rank approximation
k = 1
Mk = U[:, :k] @ np.diag(s[:k]) @ Vt[:k, :]
print(f"Best rank-{k} approximation error: {s[k]:.6f}")

# Gerschgorin discs plot
fig, ax = plt.subplots(1, 1, figsize=(8, 4))
for i in range(A.shape[0]):
    center = A[i, i]
    radius = sum(abs(A[i, j]) for j in range(A.shape[1]) if j != i)
    circle = plt.Circle((center, 0), radius, fill=True,
                         alpha=0.2, edgecolor='blue', linewidth=2)
    ax.add_patch(circle)
    ax.plot(center, 0, 'b.', markersize=8)

eigvals_exact = np.linalg.eigvalsh(A)
ax.plot(eigvals_exact, np.zeros_like(eigvals_exact), 'k*',
        markersize=12, label='Eigenvalues')
ax.set_xlim(0, 7)
ax.set_ylim(-2.5, 2.5)
ax.set_aspect('equal')
ax.axhline(y=0, color='gray', linewidth=0.5)
ax.set_xlabel('Re')
ax.set_ylabel('Im')
ax.set_title('Gerschgorin discs')
ax.legend()
plt.tight_layout()
plt.savefig('gerschgorin.pdf')
plt.show()
\end{minted}
\end{codeblock}

\begin{codeblock}[title=\textbf{Julia --- Eigenvalue algorithms}]
\begin{minted}{julia}
using LinearAlgebra, Printf

function power_iteration(A; v0=nothing, tol=1e-12, maxiter=1000)
    n = size(A, 1)
    v = isnothing(v0) ? randn(n) : copy(v0)
    v ./= norm(v)
    lam = 0.0
    for k in 1:maxiter
        w = A * v
        lam = dot(v, w)
        v = w / norm(w)
        if norm(A * v - lam * v) < tol
            @printf("Converged in %d iterations\n", k)
            break
        end
    end
    return lam, v
end

function inverse_iteration(A, sigma; tol=1e-12, maxiter=1000)
    n = size(A, 1)
    v = randn(n)
    v ./= norm(v)
    F = lu(A - sigma * I)
    lam = 0.0
    for k in 1:maxiter
        w = F \ v
        lam = dot(v, A * v)
        v = w / norm(w)
        if norm(A * v - lam * v) < tol
            @printf("Converged in %d iterations\n", k)
            break
        end
    end
    return lam, v
end

# Example
A = [4.0 1.0 0.0;
     1.0 3.0 1.0;
     0.0 1.0 2.0]

println("=== Power iteration ===")
lam, v = power_iteration(A)
@printf("Dominant eigenvalue: %.10f\n", lam)

println("\n=== Inverse iteration (sigma=1.0) ===")
lam2, v2 = inverse_iteration(A, 1.0)
@printf("Nearest eigenvalue to 1.0: %.10f\n", lam2)

println("\n=== Built-in eigvals ===")
println("Eigenvalues: ", eigvals(Symmetric(A)))

println("\n=== SVD ===")
M = [3.0 2.0 2.0; 2.0 3.0 -2.0]
F = svd(M)
println("Singular values: ", F.S)
\end{minted}
\end{codeblock}

%==============================================================================
\section{Exercises}
%==============================================================================

\begin{exercise}[$\star$]
\label{ex:ch10-gerschgorin}
For the matrix
\[
    \mathbf{A} = \begin{pmatrix} 5 & 1 & 1 \\ 0 & 4 & 1 \\ 1 & 0 & -1 \end{pmatrix},
\]
\begin{enumerate}
    \item Determine the three Gerschgorin discs.
    \item Sketch them in the complex plane.
    \item Deduce bounds on the eigenvalues. In particular, show that one eigenvalue is negative.
\end{enumerate}
\end{exercise}

\begin{exercise}[$\star$]
\label{ex:ch10-power-hand}
Apply 3 iterations of power iteration to
\[
    \mathbf{A} = \begin{pmatrix} 2 & 1 \\ 1 & 3 \end{pmatrix}, \qquad \mathbf{v}^{(0)} = \begin{pmatrix} 1 \\ 1 \end{pmatrix}.
\]
Compare the Rayleigh quotient at each step with the exact dominant eigenvalue.
\end{exercise}

\begin{exercise}[$\star\star$]
\label{ex:ch10-inverse-shift}
For the matrix in Exercise~\ref{ex:ch10-power-hand}:
\begin{enumerate}
    \item Apply inverse iteration with shift $\sigma = 1.5$ to find the smallest eigenvalue.
    \item Compare the convergence rate with unshifted power iteration.
\end{enumerate}
\end{exercise}

\begin{exercise}[$\star\star$]
\label{ex:ch10-svd-properties}
Let $\mathbf{A} = \begin{pmatrix} 1 & 2 \\ 3 & 4 \\ 5 & 6 \end{pmatrix}$.
\begin{enumerate}
    \item Compute $\mathbf{A}^\top\mathbf{A}$ and its eigenvalues.
    \item Deduce the singular values of $\mathbf{A}$.
    \item Verify with \texttt{numpy.linalg.svd}.
    \item Compute the best rank-1 approximation and its error.
\end{enumerate}
\end{exercise}

\begin{exercise}[$\star\star$]
\label{ex:ch10-qr-convergence}
Implement the basic QR algorithm and apply it to the matrix
\[
    \mathbf{A} = \begin{pmatrix} 4 & 1 & 0 & 0 \\ 1 & 4 & 1 & 0 \\ 0 & 1 & 4 & 1 \\ 0 & 0 & 1 & 4 \end{pmatrix}.
\]
\begin{enumerate}
    \item Track $\max_{i>j} |a_{ij}^{(k)}|$ (sub-diagonal elements) as a function of iteration $k$.
    \item How many iterations are needed for convergence to $10^{-12}$?
    \item Add the Wilkinson shift and compare.
\end{enumerate}
\end{exercise}

\begin{exercise}[$\star\star\star$]
\label{ex:ch10-image-compression}
\textbf{(Project: Image compression via SVD)}
\begin{enumerate}
    \item Load a grayscale image as a matrix $\mathbf{A} \in \R^{m \times n}$.
    \item Compute its SVD: $\mathbf{A} = \mathbf{U}\boldsymbol{\Sigma}\mathbf{V}^\top$.
    \item Plot the singular values (in log scale). How fast do they decay?
    \item Reconstruct the image using rank-$k$ approximations for $k = 1, 5, 10, 20, 50, 100$.
    \item For each $k$, compute the compression ratio: $\frac{k(m+n+1)}{mn}$.
    \item Plot the relative error $\|\mathbf{A} - \mathbf{A}_k\|_F / \|\mathbf{A}\|_F$ as a function of $k$.
\end{enumerate}
\end{exercise}

\begin{exercise}[$\star\star\star$]
\label{ex:ch10-hessenberg}
Implement Hessenberg reduction using Householder reflections:
\begin{enumerate}
    \item For a general $n \times n$ matrix, compute $\mathbf{H} = \mathbf{Q}^\top\mathbf{A}\mathbf{Q}$.
    \item Verify that $\mathbf{H}$ is upper Hessenberg and $\sigma(\mathbf{H}) = \sigma(\mathbf{A})$.
    \item Implement the QR iteration on the Hessenberg matrix and compare the cost with the basic QR algorithm.
\end{enumerate}
\end{exercise}

%==============================================================================
\section*{Summary of Key Formulas}
%==============================================================================

\begin{keyformulas}
\textbf{Gerschgorin:}
\[
    \sigma(\mathbf{A}) \subset \bigcup_{i=1}^n D_i, \qquad D_i = \{z \in \C : |z - a_{ii}| \leq R_i\}, \quad R_i = \sum_{j \neq i} |a_{ij}|
\]

\textbf{Power iteration:} converges to $\lambda_1$ at rate $|\lambda_2/\lambda_1|^k$ (eigenvector) or $|\lambda_2/\lambda_1|^{2k}$ (Rayleigh quotient).

\textbf{Inverse iteration with shift $\sigma$:} converges to the eigenvalue nearest $\sigma$, rate $\left|\frac{\lambda_j - \sigma}{\lambda_k - \sigma}\right|^k$.

\textbf{QR algorithm:}
\[
    \mathbf{A}_k = \mathbf{Q}_k\mathbf{R}_k, \quad \mathbf{A}_{k+1} = \mathbf{R}_k\mathbf{Q}_k \quad \to \quad \text{upper triangular (eigenvalues on diagonal)}
\]

\textbf{Wilkinson shift:} $\sigma = a_{nn} + \delta - \operatorname{sign}(\delta)\sqrt{\delta^2 + a_{n,n-1}^2}$, $\delta = (a_{n-1,n-1} - a_{nn})/2$.

\textbf{SVD:} $\mathbf{A} = \mathbf{U}\boldsymbol{\Sigma}\mathbf{V}^\top$, \quad $\|\mathbf{A}\|_2 = \sigma_1$, \quad $\cond_2(\mathbf{A}) = \sigma_1/\sigma_n$.

\textbf{Eckart--Young:} Best rank-$k$ approximation is truncated SVD, error $= \sigma_{k+1}$ (2-norm).
\end{keyformulas}


% ============================================================
\backmatter
\appendix
\chapter{Formula Sheet}
\label{app:formulas}

\section{Vector and Matrix Norms}
\begin{align}
\|x\|_1 &= \sum|x_i|, \quad \|x\|_2=\sqrt{\sum x_i^2}, \quad \|x\|_\infty=\max|x_i| \\
\|A\|_1 &= \max_j\sum_i|a_{ij}|, \quad \|A\|_\infty=\max_i\sum_j|a_{ij}|, \quad \|A\|_2=\sigma_{\max}(A)
\end{align}

\section{Errors}
\begin{align}
\text{Absolute error} &: |x-\tilde{x}| \\
\text{Relative error} &: \frac{|x-\tilde{x}|}{|x|} \\
\text{Condition number} &: \cond(A)=\|A\|\cdot\|A^{-1}\|
\end{align}

\section{Quadrature Rules}
\begin{align}
\text{Trapezoidal} &: \frac{h}{2}[f(a)+2\sum f(x_i)+f(b)], \quad E=\BigO(h^2) \\
\text{Simpson} &: \frac{h}{3}[f(a)+4\sum f_{\text{odd}}+2\sum f_{\text{even}}+f(b)], \quad E=\BigO(h^4)
\end{align}

\section{ODE Methods}
\begin{align}
\text{Explicit Euler} &: y_{n+1}=y_n+hf(t_n,y_n), \quad \BigO(h) \\
\text{RK4} &: y_{n+1}=y_n+\tfrac{h}{6}(k_1+2k_2+2k_3+k_4), \quad \BigO(h^4)
\end{align}

% --- Bibliography ---
\begin{thebibliography}{9}
\bibitem{quarteroni} \textsc{Quarteroni}, A., Sacco, R. \& Saleri, F. (2007). \emph{Numerical Mathematics}. 2nd ed., Springer.
\bibitem{suli} \textsc{S\"uli}, E. \& Mayers, D. (2003). \emph{An Introduction to Numerical Analysis}. Cambridge University Press.
\bibitem{trefethen} \textsc{Trefethen}, L.N. \& Bau, D. (1997). \emph{Numerical Linear Algebra}. SIAM.
\bibitem{higham} \textsc{Higham}, N.J. (2002). \emph{Accuracy and Stability of Numerical Algorithms}. 2nd ed., SIAM.
\bibitem{demailly} \textsc{Demailly}, J.-P. (2006). \emph{Analyse Num\'erique et \'Equations Diff\'erentielles}. EDP Sciences.
\end{thebibliography}

\printindex
\end{document}
